% concatenated from VanishingSlip-7n.tex, all-def.tex
\documentclass[pagesize,pdftex]{amsart}
%%%%%%%%%%%%%%%%%%%%%%%%%%%%%%%%%%%%%%%%%%%%%%%%%%%%%%%%%%%%%%%%%%%%%%%%%%%%%%%%%%%%%%%%%%%%%%%%%%%%%%%%%%%%%%%%%%%%%%%%%%%%%%%%%%%%%%%%%%%%%%%%%%%%%%%%%%%%%%%%%%%%%%%%%%%%%%%%%%%%%%%%%%%%%%%%%%%%%%%%%%%%%%%%%%%%%%%%%%%%%%%%%%%%%%%%%%%%%%%%%%%%%%%%%%%%
\usepackage{amssymb,amsfonts,amsmath,amsthm}
\usepackage{ulem} % remove if there are problems
\usepackage{graphicx}
%\usepackage[color]{showkeys}

\usepackage{lipsum}

\usepackage{autonum}

\usepackage{enumitem}

% ADD THE FOLLOWING COUPLE LINES INTO YOUR PREAMBLE
\let\OLDthebibliography\thebibliography
\renewcommand\thebibliography[1]{
  \OLDthebibliography{#1}
  \setlength{\parskip}{3pt}
  \setlength{\itemsep}{0pt plus 0.3ex}
}

\setcounter{MaxMatrixCols}{10}

\newtheorem{theorem}{Theorem}[section]
\newtheorem{defn}[theorem]{Definition}
\newtheorem{lem}[theorem]{Lemma}
\newtheorem{lemma}[theorem]{Lemma}
\newtheorem{prop}[theorem]{Proposition}
\newtheorem{definition}[theorem]{Definition}
\newtheorem{proposition}[theorem]{Proposition}
\newtheorem{corollary}[theorem]{Corollary}
\newtheorem{conjecture}[theorem]{Conjecture}
\newtheorem{remark}[theorem]{Remark}
\newtheorem{assumption}[theorem]{Assumption}
\newtheorem{example}[theorem]{Example}

%\input{tex/tcilatex}
\parindent=0cm

%2025.02

\newcommand{\CUS}{\Subset}

%%%%%%%%%%%%%%%%%%%%%%%%%%%%%%%%%%%%%%%%%%%%%%%%%%%%%%%%%%%%%%%%%%%%%%%%%%%
%%%%% Environments
%%%%%%%%%%%%%%%%%%%%%%%%%%%%%%%%%%%%%%%%%%%%%%%%%%%%%%%%%%%%%%%%%%%%%%%%%%%

%\newtheorem{theorem}{Theorem}[section]
%\newtheorem{definition}[theorem]{Definition}
%\newtheorem{proposition}[theorem]{Proposition}
%\newtheorem{remark}[theorem]{Remark}
%\newtheorem{lemma}[theorem]{Lemma}
%\newtheorem{corollary}[theorem]{Corollary}

\newenvironment{bb}{ \begin{beamerboxesrounded}{} }{ \end{beamerboxesrounded} }
% new itemize environment with itemsep parameter
\newenvironment{myitemize}[1][0]
  { \begin{itemize}
    \addtolength{\itemsep}{#1\baselineskip} % spacing between items
    \addtolength{\baselineskip}{#1\baselineskip} } % spacing between lines
  { \end{itemize} }
\newenvironment{myenumerate}[1][0]
  { \begin{enumerate}
    \addtolength{\itemsep}{#1\baselineskip} % spacing between items
    \addtolength{\baselineskip}{#1\baselineskip} } % spacing between lines
  { \end{enumerate} }


%\newcolumntype{x}[1]{%
%>{\raggedleft\hspace{0pt}}p{#1}}%

\newcommand{\RA}[0]{\Rightarrow} %
\newcommand{\LA}[0]{\Leftarrow} %
\newcommand{\LRA}[0]{\Leftrightarrow} %
\newcommand{\BRA}[0]{\Longrightarrow} %
\newcommand{\BLA}[0]{\Longleftarrow} %
\newcommand{\BLRA}[0]{\Longleftrightarrow} %
\newcommand{\toto}[0]{\rightrightarrows} %
\newcommand{\OL}[0]{\overline} %
\newcommand{\DBLRA}[0]{\quad :\Longleftrightarrow \quad} % 

%COMMAND SHORTCUTS
\newcommand{\qqquad}[0]{\qquad\qquad}
\newcommand{\qqqquad}[0]{\qqquad\qqquad}
\newcommand{\BS}[0]{\backslash}
\newcommand{\nt}[0]{\notag}
\newcommand{\FA}[1]{\ensuremath{\forall #1}} 
\newcommand{\FU}[1]{\text{for $#1$}}
\newcommand{\FL}[1]{\text{if $#1$}}
\newcommand{\ddx}[0]{{\frac d{dx}}}
\newcommand{\ddxk}[1]{{\frac {d^{#1}}{dx^{#1}}}}
%\newcommand{\k}[1]{{\frac {d^{#1}}{dt^{#1}}}}
\newcommand{\dde}{{\frac {d}{d\eps}}}
\newcommand{\HW}[1]{} %{\paragraph{Vorschlag für Uebungsaufgabe.}}

%%%%%%%%%%%%%%%%%%%%%%%%%%%%%%%%%%%%%%%%%%%%%%%%%%%%%%%%%%%%%%%%%%%%%%%%%%%
%%%%% kurz
%%%%%%%%%%%%%%%%%%%%%%%%%%%%%%%%%%%%%%%%%%%%%%%%%%%%%%%%%%%%%%%%%%%%%%%%%%%

\renewcommand{\t}{\tilde}
\renewcommand{\tilde}{\widetilde}
\newcommand{\p}{\partial}
\renewcommand{\d}{\delta}
\newcommand{\e}{\eps}
\newcommand{\SUS}{\subset}
\newcommand{\Rn}{\R^n}

\newenvironment{TC} {\left \{\begin{array}{ll}} {\end{array} \right.}

\newcounter{pcounter} 
\newenvironment{penum}{
  \begin{list}{(\roman{pcounter}):}{\usecounter{pcounter}}
  }{
  \end{list}
} 


\newcommand{\DS}{\displaystyle}

%%%%%%%%%%%%%%%%%%%%%%%%%%%%%%%%%%%%%%%%%%%%%%%%%%%%%%%%%%%%%%%%%%%%%%%%%%%
%%%%% mathcal
%%%%%%%%%%%%%%%%%%%%%%%%%%%%%%%%%%%%%%%%%%%%%%%%%%%%%%%%%%%%%%%%%%%%%%%%%%%

\renewcommand{\AA}{{\mathcal A}}
\newcommand{\BB}{{\mathcal B}}
\newcommand{\CC}{{\mathcal C}}
\newcommand{\DD}{{\mathcal D}}
\newcommand{\EE}{{\mathcal E}}
\newcommand{\FF}{{\mathcal F}}
\newcommand{\GG}{{\mathcal G}}
\newcommand{\HH}{{\mathcal H}}
\newcommand{\II}{{\mathcal I}}
\newcommand{\JJ}{{\mathcal J}}
\newcommand{\KK}{{\mathcal K}}
\newcommand{\LL}{{\mathcal L}}
\newcommand{\MM}{{\mathcal M}}
\newcommand{\NN}{{\mathcal N}}
\newcommand{\OO}{{\mathcal O}}
\newcommand{\PP}{{\mathcal P}}
\newcommand{\QQ}{{\mathcal Q}}
\newcommand{\RR}{{\mathcal R}}
\renewcommand{\SS}{{\mathcal S}}
\newcommand{\TT}{{\mathcal T}}
\newcommand{\UU}{{\mathcal U}}
\newcommand{\VV}{{\mathcal V}}
\newcommand{\WW}{{\mathcal W}}
\newcommand{\XX}{{\mathcal X}}
\newcommand{\YY}{{\mathcal Y}}
\newcommand{\ZZ}{{\mathcal Z}}


%%%%%%%%%%%%%%%%%%%%%%%%%%%%%%%%%%%%%%%%%%%%%%%%%%%%%%%%%%%%%%%%%%%%%%%%%%%
%%%%% mathhbb und mathbf
%%%%%%%%%%%%%%%%%%%%%%%%%%%%%%%%%%%%%%%%%%%%%%%%%%%%%%%%%%%%%%%%%%%%%%%%%%%

\newcommand{\R}{\mathbb{R}}
\newcommand{\N}{\mathbb{N}}
\newcommand{\Z}{\mathbb{Z}}
\newcommand{\C}{\mathbb{C}}
\newcommand{\T}{\mathbb{T}}
\renewcommand{\S}{\mathbb{S}}

\newcommand{\qqq}{{\mathbf q}}
\newcommand{\xxx}{{\mathbf x}}
\newcommand{\yyy}{{\mathbf y}}



%%%%%%%%%%%%%%%%%%%%%%%%%%%%%%%%%%%%%%%%%%%%%%%%%%%%%%%%%%%%%%%%%%%%%%%%%%%
%%%%% Greek symbols
%%%%%%%%%%%%%%%%%%%%%%%%%%%%%%%%%%%%%%%%%%%%%%%%%%%%%%%%%%%%%%%%%%%%%%%%%%%

\newcommand{\alp}{\alpha}
\newcommand{\bet}{\beta}
\newcommand{\gam}{\gamma}
\newcommand{\eps}{\epsilon}
\newcommand{\zet}{\zeta}
\newcommand{\kap}{\kappa}
\newcommand{\lam}{\lambda}
\renewcommand{\phi}{\varphi}
\newcommand{\ome}{\omega}
\newcommand{\sig}{\sigma}
\newcommand{\ang}{\theta}


\newcommand{\Gam}{\Gamma}
\newcommand{\Lam}{\Lambda}
\newcommand{\Sig}{\Sigma}
\newcommand{\Ome}{\Omega}

\newcommand{\pOme}{\partial \Omega}



%%%%%%%%%%%%%%%%%%%%%%%%%%%%%%%%%%%%%%%%%%%%%%%%%%%%%%%%%%%%%%%%%%%%%%%%%%%
%%%%% Math Operators
%%%%%%%%%%%%%%%%%%%%%%%%%%%%%%%%%%%%%%%%%%%%%%%%%%%%%%%%%%%%%%%%%%%%%%%%%%%

\renewcommand{\Re}{\textrm{Re }} 
\renewcommand{\Im}{\textrm{Im }} 
\renewcommand{\hat}{\widehat}
\newcommand{\nin}{\not\in}

\DeclareMathOperator*{\argmax}{argmax} 
\DeclareMathOperator{\argmin}{argmin} 
\DeclareMathOperator{\arsinh}{arsinh}
\DeclareMathOperator*{\cof}{cof} 
\DeclareMathOperator*{\DIV}{div} 
\DeclareMathOperator*{\diam}{diam}
\DeclareMathOperator*{\diag}{diag}
\DeclareMathOperator*{\dist}{dist} 
\DeclareMathOperator{\id}{Id}
\DeclareMathOperator{\Id}{Id}
\DeclareMathOperator*{\dv}{\ \nabla \- \cdot \-}
\DeclareMathOperator*{\grad}{grad} 
\DeclareMathOperator*{\graph}{graph}
\DeclareMathOperator*{\Hess}{Hess}
\DeclareMathOperator*{\kernel}{kern}
\DeclareMathOperator*{\length}{length}
\DeclareMathOperator{\Lip}{Lip}
\DeclareMathOperator*{\Mat}{Mat} 
\DeclareMathOperator*{\osc}{osc} 
\DeclareMathOperator*{\Res}{Res} 
\DeclareMathOperator*{\subgraph}{subgraph}
\DeclareMathOperator{\sign}{sign}
\DeclareMathOperator*{\spt}{supp}
\DeclareMathOperator*{\spann}{span}
\DeclareMathOperator*{\SO}{SO} 
\DeclareMathOperator*{\Skew}{Skew} 
\DeclareMathOperator*{\Sym}{Sym} 

\def \wto{\rightharpoonup}
%\renewcommand{\Re}{\mathfrak{Re}}
%\renewcommand{\Im}{\mathfrak{Im}}
\newcommand{\da}{\partial}
\newcommand{\dd}{\, {\rm d}}
\newcommand{\ddt}{\, \frac{d}{dt}}
\renewcommand{\iint}{\int\!\!\!\!\int}
\newcommand{\gauss}[1]{\lfloor #1 \rfloor}                        % abrunden
\newcommand{\avg}[1]{\langle #1 \rangle}
\newcommand{\avgL}[2]{\langle #1 \rangle_{#2}}


% Mean value
\def\Xint#1{\mathchoice
{\XXint\displaystyle\textstyle{#1}}%
{\XXint\textstyle\scriptstyle{#1}}%
{\XXint\scriptstyle\scriptscriptstyle{#1}}%
{\XXint\scriptscriptstyle\scriptscriptstyle{#1}}%
\!\int}
\def\XXint#1#2#3{{\setbox0=\hbox{$#1{#2#3}{\int}$}
\vcenter{\hbox{$#2#3$}}\kern-.5\wd0}}
\def\ddashint{\Xint=}
\def\dashint{\Xint-}

% \upref, \upeq

\newcommand{\upref}[2]{\hspace{-0.8ex}\stackrel{\eqref{#1}}{#2}} % standard
\newcommand{\xupref}[2]{\hspace{-0.3ex}\stackrel{\eqref{#1}}{#2}} % for (x)
\newcommand{\upupref}[3]{\hspace{-3ex}\stackrel{\eqref{#1},\eqref{#2}}{#3}}
\newcommand{\upeq}[2]{\stackrel{#1}{#2}}
\newcommand{\uptext}[2]{\stackrel{\text{#1}}{#2}}

\newcommand{\lupref}[2]{\hspace{0ex} \stackrel{\eqref{#1}}{#2}} % in line
\newcommand{\lupupref}[3]{\hspace{0ex}\stackrel{\eqref{#1},\eqref{#2}}{#3}}
\newcommand{\luptext}[2]{\hspace{0ex} \stackrel{\text{#1}}{#2}} % in line

%%%%%%%%%%%%%%%%%%%%%%%%%%%%%%%%%%%%%%%%%%%%%%%%%%%%%%%%%%%%%%%%%%%%%%%%%%%
% NEU Farbe
%%%%%%%%%%%%%%%%%%%%%%%%%%%%%%%%%%%%%%%%%%%%%%%%%%%%%%%%%%%%%%%%%%%%%%%%%%%

\usepackage{color}
% blue
\definecolor{verylightblue}{rgb}{0.95, 0.95, 0.95}  %{0.9 0.95 1}
\definecolor{lightblue}{rgb}{0.7, 0.7, 1}

% yellow
\definecolor{eqyellow}{rgb}{0.9375,0.8984,0.5469}
\definecolor{subeqyellow}{rgb}{1,0.9373,0.8353}

%cyan
\definecolor{darkcyan}{rgb}{0.0, 0.45, 0.95} % modified from {0.0, 0.55, 0.55}
\newcommand{\darkcyan}{\color{darkcyan}}

% green
\definecolor{mygreen}{rgb}{0.3, 0.6, 0.3} %{0.9 1 0.9}
\definecolor{verylightgreen}{rgb}{0.95, 0.95, 0.95} %{0.9 1 0.9}
\definecolor{verydarkgreen}{rgb}{0, 0.5, 0}
\definecolor{darkgreen}{rgb}{0.1, 0.55, 0.1}  %{0 0.5 0}
\definecolor{mydarkgreen}{rgb}{0, 0.5, 0} %{0.9 1 0.9}

%brown
\definecolor{mybrown}{rgb}{0.85, 0.4, 0.3}
\definecolor{verylightbrown}{rgb}{0.98, 0.72, 0.58}
\definecolor{verydarkbrown}{rgb}{0.44, 0.26, 0.26}

% orange
\definecolor{orange}{rgb}{1, 0.5, 0}
\newcommand{\orange}[1]{{\color{orange}{#1}}}
\definecolor{BurntOrange}{rgb}{1,0.356,0}
\newcommand{\BurntOrange}{\color{BurntOrange}}

%red
\definecolor{mydarkred}{rgb}{1,0.086,0.255}
\newcommand{\mydarkred}{\color{mydarkred}}
\definecolor{RoseVYDP}{rgb}{0.84,0.086,0.255}
\newcommand{\RoseVYDP}{\color{RoseVYDP}}
\definecolor{dgreen}{rgb}{0, 0.8, 0.5}     
\newcommand{\dgreen}{\color{dgreen}}
\definecolor{CanaryBRT}{rgb}{1,0.76,0.26}
\newcommand{\CanaryBRT}{\color{CanaryBRT}}

% other
\definecolor{cyan}{rgb}{0, 1, 1}
\definecolor{verylightgray}{rgb}{0.95, 0.95, 0.95}
\definecolor{lightgray}{rgb}{0.8, 0.8, 0.8}
\definecolor{verylightred}{rgb}{1, 0.8, 0.78}
\definecolor{verylightyellow}{rgb}{0.99, 0.98, 0.5}


\newcommand{\black}{\color{black}}
\newcommand{\white}{\color{white}}
\newcommand{\mygreen}{\color{mygreen}}
\newcommand{\mydarkgreen}{\color{mydarkgreen}}
\newcommand{\mybrown}{\color{mybrown}}
\newcommand{\blue}[1]{{\color{blue}{#1}}}
\newcommand{\red}[1]{{\color{red}{#1}}}
\newcommand{\green}[1]{{\color{green}{#1}}}
\newcommand{\darkgreen}[1]{{\color{darkgreen}{#1}}}
\newcommand{\cyan}[1]{{\color{cyan}{#1}}}
\newcommand{\lightgray}[1]{{\color{lightgray}{#1}}}

%%%%%%%%%%%%%%%%%%%%%%%%%%%%%%%%%%%%%%%%%%%%%%%%%%%%%%%%%%%%%%%%%%%%%%%%%%%
%%%%% edit
%%%%%%%%%%%%%%%%%%%%%%%%%%%%%%%%%%%%%%%%%%%%%%%%%%%%%%%%%%%%%%%%%%%%%%%%%%%

\newcommand{\edit}[1]{{\color{red}{#1}}}
\newcommand{\short}[1]{{\color{orange}{#1}}}
\newcommand{\ignore}[1]{{}}

%%%%%%%%%%%%%%%%%%%%%%%%%%%%%%%%%%%%%%%%%%%%%%%%%%%%%%%%%%%%%%%%%%%%%%%%%%%
%%%%% flush right without line break  "\signed"
%%%%%%%%%%%%%%%%%%%%%%%%%%%%%%%%%%%%%%%%%%%%%%%%%%%%%%%%%%%%%%%%%%%%%%%%%%%


\parindent0em

\newcommand{\signed}[1]{%
  \unskip\hspace*{1em plus 1fill}%
  \nolinebreak[3]\hspace*{\fill}\mbox{\upshape #1}
}

\newenvironment{xquote}{%
  \begin{list}{}{%
    \rightmargin
    \leftmargin
  }
  \itshape
  \item[]
}
{\end{list}}


%%%%%%%%%%%%%%%%%%%%%%%%%%%%%%%%%%%%%%%%%%%%%%%%%%%%%%%%%%%%%%%%%%%%%%%%%%%
%%%%% Standard norms
%%%%%%%%%%%%%%%%%%%%%%%%%%%%%%%%%%%%%%%%%%%%%%%%%%%%%%%%%%%%%%%%%%%%%%%%%%%

%\newcommand{\NN}[1]{\|#1\|}
\newcommand{\NNN}[2]{\|#1\|_{#2}}
\newcommand{\MMM}[2]{[#1]_{#2}}
\newcommand{\skp}[2]{(#1, #2)}
\newcommand{\skpL}[3]{\skp{#1}{#2}_{#3}}
\newcommand{\SKP}[2]{\langle #1, #2 \rangle}
\newcommand{\SKPL}[3]{\SKP{#1}{#2}_{#3}}


\newcommand{\NT}[1]{\|#1\|_{L^2}}
\newcommand{\NTL}[2]{\|#1\|_{L^2({#2})}}
\newcommand{\NI}[1]{\|#1\|_{L^\infty}}
\newcommand{\NIL}[2]{\|#1\|_{L^\infty({#2})}}
\newcommand{\NP}[2]{\|#1\|_{L^{#2}}}
\newcommand{\NPL}[3]{\|#1\|_{L^{#2}(#3)}}
\newcommand{\NBV}[1]{\|\nabla #1\|_{\mathcal M}}
\newcommand{\NBVL}[2]{\|#1\|_{BV({#2})}}
\newcommand{\cci}{C_c^{\infty}(\R)}
\newcommand{\cciL}[1]{C_c^{\infty}(#1)}

\makeatletter
\newcommand{\VEC}[2][r]{% originaler Name Spvek
  \gdef\@VORNE{1}
  \left(\hskip-\arraycolsep%
    \begin{array}{#1}\vekSp@lten{#2}\end{array}%
  \hskip-\arraycolsep\right)}
%%% Vektoren 
%%% http://tex.stackexchange.com/questions/2705/typesetting-column-vector
\def\vekSp@lten#1{\xvekSp@lten#1;vekL@stLine;}
\def\vekL@stLine{vekL@stLine}
\def\xvekSp@lten#1;{\def\temp{#1}%
  \ifx\temp\vekL@stLine
  \else
    \ifnum\@VORNE=1\gdef\@VORNE{0}
    \else\@arraycr\fi%
    #1%
    \expandafter\xvekSp@lten
  \fi}
\makeatother





%%%%%%%%%%%%%%%%%%%%%%%%%%%%%%%%%%%%%%%%%%%%%%%%%%%%%%%%%%%%%%%%%%%%%%%%%%%
%%%%% H-norms
%%%%%%%%%%%%%%%%%%%%%%%%%%%%%%%%%%%%%%%%%%%%%%%%%%%%%%%%%%%%%%%%%%%%%%%%%%%

\newcommand{\HS}[1]{H_{#1}}
\newcommand{\HM}[2]{[#1]_{\HS{#2}}}
\newcommand{\HN}[2]{\|#1\|_{\HS{#2}}}
\newcommand{\HP}[3]{\langle #1, #2 \rangle_{\HS{#3}}}

\newcommand{\HSL}[2]{\HS{#1}({#2})}
\newcommand{\HML}[3]{[#1]_{\HSL{#2}{#3}}}
\newcommand{\HNL}[3]{\|#1\|_{\HSL{#2}{#3}}}
\newcommand{\HPL}[4]{\langle #1, #2 \rangle_{\HSL{#3}{#4}}}

\newcommand{\HSS}[1]{H_{#1*}}
\newcommand{\HMS}[2]{[#1]_{\HSS{#2}}}
\newcommand{\HNS}[2]{\|#1\|_{\HSS{#2}}}
\newcommand{\HPS}[3]{\langle #1, #2 \rangle_{\HSS{#3}}}

\newcommand{\HSP}[1]{H_{#1}^+}
\newcommand{\HNP}[2]{\|#1\|_{\HSP{#2}}}
\newcommand{\HSM}[1]{H_{#1}^-}
\newcommand{\HNM}[2]{\|#1\|_{\HSM{#2}}}

%\newcommand{\HSo}[1]{H_{0}^{#1}}
%\newcommand{\HMo}[2]{[#1]_{\HS{#2}}}
%\newcommand{\HNo}[2]{\|#1\|_{\HS{#2}}}

\newcommand{\HSoo}[1]{\mathring{H}_{#1}}
\newcommand{\HMoo}[2]{[#1]_{\HS{#2}}}
\newcommand{\HNoo}[2]{\|#1\|_{\HS{#2}}}

\newcommand{\HSooL}[2]{\mathring H_{#1}({#2})}
\newcommand{\HMooL}[3]{[#1]_{\HSooL{#2}{#3}}}
\newcommand{\HNooL}[3]{\|#1\|_{\HSooL{#2}{#3}}}
\newcommand{\HPooL}[4]{\langle #1, #2 \rangle_{\HSooL{#3}{#4}}}

%% time-space

% H^k

\newcommand{\THS}[1]{L^2(\HS{#1})}
\newcommand{\THM}[2]{[#1]_{\THS{#2}}}
\newcommand{\THN}[2]{\|#1\|_{\THS{#2}}}

\newcommand{\THSK}[2]{H^{#1}(\HS{#2})}
\newcommand{\THMK}[3]{[#1]_{\THSK{#2}{#3}}}
\newcommand{\THNK}[3]{\|#1\|_{\THSK{#2}{#3}}}

\newcommand{\THSS}[1]{L^2(\HSS{#1})}
\newcommand{\THMS}[2]{[#1]_{\THSS{#2}}}
\newcommand{\THNS}[2]{\|#1\|_{\THSS{#2}}}

% C^k

\newcommand{\CHS}[1]{C^0(\HS{#1})}
\newcommand{\CHM}[2]{[#1]_{\CHS{#2}}}
\newcommand{\CHN}[2]{\|#1\|_{\CHS{#2}}}

\newcommand{\CHSK}[2]{C^{#1}(\HS{#2})}
\newcommand{\CHMK}[3]{[#1]_{\CHSK{#2}{#3}}}
\newcommand{\CHNK}[3]{\|#1\|_{\CHSK{#2}{#3}}}

\newcommand{\CHSS}[1]{C^0(\HSS{#1})}
\newcommand{\CHMS}[2]{[#1]_{\CHSS{#2}}}
\newcommand{\CHNS}[2]{\|#1\|_{\CHSS{#2}}}

%% parabolic

\newcommand{\PHS}[1]{TH_{#1}}
\newcommand{\PHM}[2]{[#1]_{\dot \PHS{#2}}}
\newcommand{\PHN}[2]{\|#1\|_{\PHS{#2}}}

\newcommand{\PHSL}[2]{TH_{#1}({#2})}
\newcommand{\PHML}[3]{[#1]_{\dot \PHSL{#2}{#3}}}
\newcommand{\PHNL}[3]{\|#1\|_{\PHSL{#2}{#3}}}


%% parabolic-from 0

\newcommand{\PHSN}[1]{\PHS{#1} \cap L^2(H_\eps^0)}
\newcommand{\PHMN}[2]{[#1]_{\dot \PHSN{#2}}}
\newcommand{\PHNN}[2]{\|#1\|_{\PHSN{#2}}}
 

%%%%%%%%%%%%%%%%%%%%%%%%%%%%%%%%%%%%%%%%%%%%%%%%%%%%%%%%%%%%%%%%%%%%%%%%%%%
%%%%% X-norms
%%%%%%%%%%%%%%%%%%%%%%%%%%%%%%%%%%%%%%%%%%%%%%%%%%%%%%%%%%%%%%%%%%%%%%%%%%%

\newcommand{\XS}[1]{X_{#1}}
\newcommand{\XM}[2]{[#1]_{\XS{#2}}}
\newcommand{\XN}[2]{\|#1\|_{\XS{#2}}}
\newcommand{\XP}[3]{\langle #1, #2 \rangle_{\XS{#3}}}

\newcommand{\XSL}[2]{X_{#1}({#2})}
\newcommand{\XML}[3]{[#1]_{\XSL{#2}{#3}}}
\newcommand{\XNL}[3]{\|#1\|_{\XSL{#2}{#3}}}
\newcommand{\XPL}[4]{\langle #1, #2 \rangle_{\XSL{#3}{#4}}}




\newcommand{\XSP}[1]{X_{#1}^+}
\newcommand{\XNP}[2]{\|#1\|_{\XSP{#2}}}
\newcommand{\XSM}[1]{X_{#1}^-}
\newcommand{\XNM}[2]{\|#1\|_{\XSM{#2}}}


\newcommand{\XSoo}[1]{\mathring{X}_{#1}}
\newcommand{\XMoo}[2]{[#1]_{\XS{#2}}}
\newcommand{\XNoo}[2]{\|#1\|_{\XS{#2}}}

%% time-space

% X^k

\newcommand{\TXS}[1]{L^2(\XS{#1})}
\newcommand{\TXM}[2]{[#1]_{\TXS{#2}}}
\newcommand{\TXN}[2]{\|#1\|_{\TXS{#2}}}
\newcommand{\TXSL}[2]{L_t^2\XS{#1}({#2})}
\newcommand{\TXNL}[3]{\|#1\|_{\TXSL{#2}{#3}}}

\newcommand{\TXSK}[2]{X_{#1}(\XS{#2})}
\newcommand{\TXMK}[3]{[#1]_{\TXSK{#2}{#3}}}
\newcommand{\TXNK}[3]{\|#1\|_{\TXSK{#2}{#3}}}

\newcommand{\TXSS}[1]{L^2(\XSS{#1})}
\newcommand{\TXMS}[2]{[#1]_{\TXSS{#2}}}
\newcommand{\TXNS}[2]{\|#1\|_{\TXSS{#2}}}

% C^k

\newcommand{\CXS}[1]{C^0(\XS{#1})}
\newcommand{\CXM}[2]{[#1]_{\CXS{#2}}}
\newcommand{\CXN}[2]{\|#1\|_{\CXS{#2}}}
\newcommand{\CXSL}[2]{C_t^0\XS{#1}({#2})}
\newcommand{\CXNL}[3]{\|#1\|_{\CXSL{#2}{#3}}}

\newcommand{\CXSK}[2]{C^{#1}(\XS{#2})}
\newcommand{\CXMK}[3]{[#1]_{\CXSK{#2}{#3}}}
\newcommand{\CXNK}[3]{\|#1\|_{\CXSK{#2}{#3}}}

\newcommand{\CXSS}[1]{C^0(\XSS{#1})}
\newcommand{\CXMS}[2]{[#1]_{\CXSS{#2}}}
\newcommand{\CXNS}[2]{\|#1\|_{\CXSS{#2}}}

%% parabolic

\newcommand{\PXS}[1]{TX_{#1}}
\newcommand{\PXM}[2]{[#1]_{\dot \PXS{#2}}}
\newcommand{\PXN}[2]{\|#1\|_{\PXS{#2}}}

\newcommand{\PXSL}[2]{\PXS{#1}({#2})}
\newcommand{\PXML}[3]{[#1]_{\dot \PXSL{#2}{#3}}}
\newcommand{\PXNL}[3]{\|#1\|_{\PXSL{#2}{#3}}}


%% parabolic-from 0

\newcommand{\PXSN}[1]{\PXS{#1} \cap L^2(X_\eps^0)}
\newcommand{\PXMN}[2]{[#1]_{\dot \PXSN{#2}}}
\newcommand{\PXNN}[2]{\|#1\|_{\PXSN{#2}}}

%%%%%%%%%%%%%%%%%%%%%%%%%%%%%%%%%%%%%%%%%%%%%%%%%%%%%%%%%%%%%%%%%%%%%%%%%%%
%%%%% Y-norms
%%%%%%%%%%%%%%%%%%%%%%%%%%%%%%%%%%%%%%%%%%%%%%%%%%%%%%%%%%%%%%%%%%%%%%%%%%%

\newcommand{\YS}[1]{Y_{#1}}
\newcommand{\YM}[2]{[#1]_{\YS{#2}}}
\newcommand{\YN}[2]{\|#1\|_{\YS{#2}}}
\newcommand{\YP}[3]{\langle #1, #2 \rangle_{\YS{#3}}}

\newcommand{\YSL}[2]{\YS{#1}({#2})}
\newcommand{\YML}[3]{[#1]_{\YSL{#2}{#3}}}
\newcommand{\YNL}[3]{\|#1\|_{\YSL{#2}{#3}}}
\newcommand{\YPL}[4]{\langle #1, #2 \rangle_{\YSL{#3}{#4}}}

\newcommand{\YSS}[1]{Y_{#1*}}
\newcommand{\YMS}[2]{[#1]_{\YSS{#2}}}
\newcommand{\YNS}[2]{\|#1\|_{\YSS{#2}}}
\newcommand{\YPS}[3]{\langle #1, #2 \rangle_{\YSS{#3}}}


\newcommand{\YSP}[1]{Y_{#1}^+}
\newcommand{\YNP}[2]{\|#1\|_{\YSP{#2}}}
\newcommand{\YSM}[1]{Y_{#1}^-}
\newcommand{\YNM}[2]{\|#1\|_{\YSM{#2}}}

% homoegenous bc

\newcommand{\YSoo}[1]{\mathring{Y}_{#1}}
\newcommand{\YMoo}[2]{[#1]_{\YS{#2}}}
\newcommand{\YNoo}[2]{\|#1\|_{\YS{#2}}}

\newcommand{\PYSL}[2]{TY_{#1}({#2})}
\newcommand{\PYML}[3]{[#1]_{\dot \PYSL{#2}{#3}}}
\newcommand{\PYNL}[3]{\|#1\|_{\PYSL{#2}{#3}}}

%% time-space

% Y^k

\newcommand{\TYS}[1]{L^2(\YS{#1})}
\newcommand{\TYM}[2]{[#1]_{\TYS{#2}}}
\newcommand{\TYN}[2]{\|#1\|_{\TYS{#2}}}

\newcommand{\TYSK}[2]{Y_{#1}(\YS{#2})}
\newcommand{\TYMK}[3]{[#1]_{\TYSK{#2}{#3}}}
\newcommand{\TYNK}[3]{\|#1\|_{\TYSK{#2}{#3}}}

\newcommand{\TYSS}[1]{L^2(\YSS{#1})}
\newcommand{\TYMS}[2]{[#1]_{\TYSS{#2}}}
\newcommand{\TYNS}[2]{\|#1\|_{\TYSS{#2}}}

% C^k

\newcommand{\CYS}[1]{C^0(\YS{#1})}
\newcommand{\CYM}[2]{[#1]_{\CYS{#2}}}
\newcommand{\CYN}[2]{\|#1\|_{\CYS{#2}}}

\newcommand{\CYSK}[2]{C^{#1}(\YS{#2})}
\newcommand{\CYMK}[3]{[#1]_{\CYSK{#2}{#3}}}
\newcommand{\CYNK}[3]{\|#1\|_{\CYSK{#2}{#3}}}

\newcommand{\CYSS}[1]{C^0(\YSS{#1})}
\newcommand{\CYMS}[2]{[#1]_{\CYSS{#2}}}
\newcommand{\CYNS}[2]{\|#1\|_{\CYSS{#2}}}

%% parabolic

\newcommand{\PYS}[1]{TY_{#1}}
\newcommand{\PYM}[2]{[#1]_{\dot \PYS{#2}}}
\newcommand{\PYN}[2]{\|#1\|_{\PYS{#2}}}


%% parabolic-from 0

\newcommand{\PYSN}[1]{\PYS{#1} \cap L^2(Y_\eps^0)}
\newcommand{\PYMN}[2]{[#1]_{\dot \PYSN{#2}}}
\newcommand{\PYNN}[2]{\|#1\|_{\PYSN{#2}}}


%%%%%%%%%%%%%%%%%%%%%%%%%%%%%%%%%%%%%%%%%%%%%%%%%%%%%%%%%%%%%%%%%%%%%%%%%%%
%%%%% Z-norms
%%%%%%%%%%%%%%%%%%%%%%%%%%%%%%%%%%%%%%%%%%%%%%%%%%%%%%%%%%%%%%%%%%%%%%%%%%%



\newcommand{\ZS}[1]{Z^{#1}}
\newcommand{\ZM}[2]{[#1]_{\ZS{#2}}}
\newcommand{\ZN}[2]{\|#1\|_{\ZS{#2}}}
\newcommand{\ZP}[3]{\langle #1, #2 \rangle_{\ZS{#3}}}

\newcommand{\ZSL}[2]{Z^{#1}({#2})}
\newcommand{\ZML}[3]{[#1]_{\ZS{#2}{#3}}}
\newcommand{\ZNL}[3]{\|#1\|_{\ZS{#2}{#3}}}
\newcommand{\ZPL}[4]{\langle #1, #2 \rangle_{\ZS{#3}{#4}}}


\newcommand{\hk}[1]{{\blue{#1}}}
\newcommand{\hkc}[1]{{\darkcyan{$\langle\hspace{-0.3ex}\langle$#1$\rangle\hspace{-0.3ex}\rangle$}}}
\newcommand{\hkrm}[1]{\hk{\xout{#1}}} %
\newcommand{\QUESTION}[1]{#1}
\newcommand{\COMMENT}[1]{#1}
\DeclareMathOperator{\supp}{supp}
\newcommand{\Ai}{{\rm Ai}}
\newcommand{\Bi}{{\rm Bi}}

\renewcommand{\theenumi}{\roman{enumi}}
\renewcommand{\labelenumi}{\textrm(\theenumi)}
\newcommand{\LAM}{{\mathfrak s}}

%\usepackage{setspace}
%\doublespacing


\title{Solutions of the thin film equation obtained in the limit of vanishing slip}
\subjclass[2020]{35Q35,35K65,35C20,35B25}
\keywords{Fluid dynamics, Degenerate Parabolic Equations, Thin--film Equation, Singular Limit}

\author{Hans Knüpfer}
\address{Institut für Angewandte Mathematik \& Interdisciplinary Center for Scientific Computing, Universität Heidelberg,  
Im Neuenheimer Feld 205, 69120 Heidelberg, Germany}
\email{knuepfer@uni-heidelberg.de}
\author{Juan J.~L. Velázquez}
\address{Institut für Angewandte Mathematik, Rheinische Friedrich-Wilhelms-Universität Bonn,  
Endenicher Allee 60, 53115 Bonn, Germany}
\email{velazquez@iam.uni-bonn.de}

\begin{document}

\begin{abstract}
  We analyze the evolution of thin liquid droplets in the lubrication
  approximation with different slip conditions at the liquid--solid
  interface. Motivated by the classical no-slip paradox which states that the
  Navier–Stokes equations with a no-slip boundary condition require unphysical
  infinite dissipation during droplet spreading, we focus on the limit of
  vanishing slip.  We show that in the no--slip limit three fundamentally
  different classes of limiting solutions are approached, each of them
  corresponding to a different scaling of the microscopic contact angle as the
  regularization parameter vanishes. These findings suggest that the thin-film
  equation with no slip supports a rich family of physically admissible
  solutions, provided one interprets the no--slip thin film equation as the
  asymptotic limit of models which regularized slip conditions. Even though the
  large apparent contact angles in some of these solutions seem incompatible
  with the lubrication approximation, a refined analysis shows that the
  underlying physical variables remain consistent with the assumptions for the
  lubrication approximation.
\end{abstract}
 
\renewcommand{\em}[1]{\textit{#1}}


\maketitle
\tableofcontents


\section{Introduction and main results}


% Some general questions
% \begin{enumerate}
% \item Mathematical accuracy.  Range of validity of $n$ in the equations
% \item Consistent notation (after many changes). In particular $=$ vs $\approx$
%   vs $+ \OO(\ldots)$ and $\eps \ll 1$ vs $\eps \to 0$.
% \item References: Did we cite all necessary literature? We can have a look
%   together at the references \cite{GiacomelliGnannPeschka-2023,BowenKing-2001}.
%   Should we remove the reference to our planned ODE paper? Even the authors are
%   not clear.
% \end{enumerate}



\subsection{Introduction}

We consider the evolution of a liquid droplet on a solid substrate. In the
classical fluid dynamical description, the evolution in the fluid is described
by the Navier--Stokes equations, while a no slip condition is assumed at the
liquid--solid interface. It has been observed that such a no--slip condition
allows for droplet spreading only under the seemingly unphysical condition of
infinite dissipation of energy, e.g. \cite{HuhScriven-1971}. This apparent
\textit{no slip paradox} has led to many discussions about possible resolutions
of this paradox. We want to discuss some related issues. As one way to
avoid the no--slip paradox different relaxed slip conditions have been proposed
where the tangential velocity, relative to the solid, at the interface is a
function of the tangential shear stress. In particular, we consider the
evolution of a thin capillary driven droplet in the so--called lubrication
approximation regime. Assuming a no slip condition, in this regime the evolution
of the droplet can be described by the thin--film equation
\begin{align} \label{tfe-noslip} 
    h_t +(h^3 h_{xxx})_x \ = \ 0 \qquad\qquad %
  \text{in $\Ome \ := \ \{h>0\}$.}
\end{align}
Here, $h$ is the droplet height and the continuity assumption $h = 0$ is assumed
at $\p \Ome$.  Assuming more general slip conditions, the lubrication
approximation leads to the family of equations
\begin{align} \label{tfe-slip} 
  h_t +((h^3+ \eps^{3-n} h^{n})h_{xxx})_x \ = \ 0\qquad\qquad %
  \text{in $\Ome \ := \ \{h>0\}$,}
\end{align}
The mobility exponent $n \in (0,3)$ models the slip condition at the
liquid--solid interface: $n = 2$ represents Navier--slip where a linear
slip/stress relationship is assumed, while $n = 1$ corresponds to Greenspan's
slip condition \cite{Greenspan-1978}. The regularization parameter $\eps > 0$,
which is assumed to be small, is called \textit{slip length} (see Section
\ref{LubApp}). We recover the equation \eqref{tfe-noslip} in the limit
$\eps \to 0$. Equation \eqref{tfe-slip} has to be complemented by a boundary
condition at $\p \Ome$. We assume the constant angle condition
\begin{align} \label{tfe-bc} %
  h_x \ = \ \ang_\eps \qquad\qquad\qquad %
    \text{at $\p \Ome \ = \ \p\{h>0\}$.}
\end{align}
We note that the boundary condition \eqref{tfe-bc} for \eqref{tfe-noslip} leads
to an ill--posed problem \cite{EggersFontelos-2025}. Throughout the paper, for
both equations \eqref{tfe-noslip} and \eqref{tfe-slip} we also use the boundary
condition
\begin{align} \label{no-massflux} %
  h^n h_{xxx} \ = \ 0 \qquad\qquad \text{for $x \in \Ome$, $x \to \p \Ome$.}
\end{align}
This condition ensures that there is no loss of mass across the contact points
$\p \Ome$. For further references we refer e.g. to
\cite{BDQ-Book,HuhScriven-1971,Hocking-1983,EggersStone-2004}, for nonlinear
slip see e.g. \cite{ThompsonTroian-1997}. We note that other regularizations of
the no--slip paradox exist which are not considered here, see
e.g. \cite{OronDavisBankoff-1997,WitelskiBernoff-1999,OronBankoff-2001,BowenKing-2001}. For
existence of solutions for \eqref{tfe-noslip} we refer e.g. to
\cite{BertozziPugh-1994,BerettaBertschDalpasso-1995,BertschDalpassoGarckeGruen-1998,DalpassoGarckeGruen-1998,Gruen-2004-1,GiacomelliKnuepferOtto-2008,GnannWisse-2022,GnannWisse-2025}. The
literature is discussed more in detail in the next section.


\medskip

More in particular, we study solutions of \eqref{tfe-noslip} which appear as a
limit of solutions for \eqref{tfe-slip} as $\eps \to 0$. Our arguments are
formal in the sense that they are based on the use of matched asymptotic
expansions.  For simplicity, we consider solutions with an isolated contact
point $s(t)$ at the left boundary of the support, i.e
$\Ome \cap B_{\d}(s(t)) = [s(t), s(t) + \d)$ for some $\d >
0$. We obtain the following three types of solutions of \eqref{tfe-noslip} and
their asymptotic behaviour near the contact point:
\begin{itemize}
\item[(a)] \label{Typea} %
  The support is constant, $s(t)= s_{0}$. Furthermore, 
\begin{align} \label{S1E4} %
  h(x,t) \ \approx \ \gam(t) \left(x-s_{0}\right) \qquad\qquad %
  \text{ as $x \to  s_0^+$}. 
\end{align}%
for some smooth function $\gam(t)$ (see the paragraph on notation below
for the meaning of the symbol ``$\approx$''). We recall the well--known result
that the support is constant for sufficiently regular solutions (Theorem
\ref{thm-nomove}).
\end{itemize}
\begin{itemize}
\item[(b)] \label{Typeb} %
  The contact point $s(t)$ is an arbitrary, monotonically increasing,
  prescribed smooth function (this corresponds to a shrinking droplet). The
  solutions have asymptotic behaviour %
  \begin{align}
    h(x,t) \ %
    \approx \ (3 \dot s(t))^{\frac 13}\left(x-s(t) \right) \Big(\ln \frac{1}{
    x-s(t)} \Big)^{\frac 13} \qquad\qquad %
    \text{ as $x \to  s(t)^+$}.  \label{def-typeb}
  \end{align}
  (with notation $\dot s = \frac d{dt} s$) We remark that the possibility of
  having travelling wave solutions with receding support, as the type (b)
  solutions was already suggested in \cite{BoattoEtal-1993} and obtained by
  matched asymptotics in \cite{BowenKing-2001}. It turns out that the same
  initial data allows for infinitely many solutions.
\end{itemize}
\begin{itemize}
\item[(c)] \label{Typec} %
  The interface $s$ is also monotonically increasing and moves in a shorter
  time scale $\tau$ than $t,$ that will be denoted as $\tau $ and the motion as
  $s\left(\tau \right)$. The mass accumulates in the moving interface
  $s(\tau)$. Therefore (with the notation $s' := \frac{d}{d\tau}s$) %
\begin{align}
  h(x,\tau) =m(\tau) \delta_{s(\tau)}(x) +h_{0}(x) \chi_{\{ x>s\left(\tau
  \right)\} }  \ \  %
  \text{with} \ m'(\tau) = h_{0}(s(\tau)) s'(\tau)
\label{def-typec}
\end{align}
For these solutions the interface moves so fast that the profile $h_{0}$ does
not have time to change. In this case all the mass of $%
h_{0}$ that is reduced due to the reduction of the support accumulates at the
interface boundary. This requires a more sophisticated analysis, see the
upcoming paper \cite{paper2}.
\end{itemize}
The above solutions of \eqref{tfe-noslip} appear as limit of solutions of
\eqref{tfe-slip} and are related to the asymptotic behaviour of the contact
angle $\ang_\eps$ in the limit $\eps \to 0$.

\medskip

In fact, a contact angle $\ang_\eps \sim 1$ of order corresponds to type (a)
solutions. A contact angle of size $\ang _{\eps } \gg |\ln \eps|^{\frac 13}$
leads to type (c) solutions. A contact angle in the intermediate regime
$\ang _{\eps } \approx \gam |\ln \eps|^{\frac 13}$ leads to type (b)
solutions. The contact point velocity in the intermediate regime is given by
\begin{align} \label{dots-formula} %
  \dot s \ %
  = \   + \lim_{\eps \to 0} \frac {\ang_\eps^3}{3|\ln \eps|} \ %
  = \ + \frac \gam 3 \ > \ 0, \qquad %
  \text{where $\ang_\eps \approx \gam (-\ln \eps)^{\frac 13}$.}
\end{align} 
We believe that type (a) solutions also appear in the remaining case where
$1 \ll \ang_\eps \ll |\ln \eps|^{\frac 13}$ but this case is not treated in this
paper. We also note that the limits above might only exist for a subsequence
$\eps_k \to 0$ where the limiting angle has a well--defined behaviour. Assuming
that \eqref{tfe-slip} is a valid physical regularization of \eqref{tfe-noslip}
this implies that the three types of solutions (a), (b) and (c) are physically
feasible as well.
\begin{remark}[Time scales] \label{rem-time}\text{} %
  The evolution exhibits two time scales. In our (non--dimensionalized) model,
  the droplet profile $h$ changes in $\OO(1)$ time scales, i.e. times scales of
  order one. The motion of the contact point, however, generally lives on a
  different time scale which depends on the value of the microscopic contact
  angle.  This second time scale is associated with large times for type (a)
  solutions, $\OO(1)$ times scales for type (b) solutions and much shorter times
  for type (c) solutions.
\end{remark}
Equations \eqref{tfe-noslip} and \eqref{tfe-slip} describe an evolution which is
governed by the dissipation of capillary energy. The energy--dissipation formula
for \eqref{tfe-slip} has the form%
\begin{align} \label{diss-intro} %
  \frac d{dt} \frac{1}{2}\int_{\Ome }h_x^{2}  %
   +  \ang_{\eps }^{2} \ dx  \ %
  = \ -\int_{\Ome }(h^3+ \eps^{3-n} h^{n})h_{xxx}^{2}\ dx,
\end{align}%
assuming that the solutions are sufficiently regular. Here, the two terms on the
left--hand side describe the (relative) interfacial energies between the three
phases air, liquid and solid. The right--hand side describes the dissipation of
energy due to viscous friction.  For type (b) solutions of \eqref{tfe-noslip} we
obtain that the rate of dissipation of energy near the contact point is
divergent for type (b) solutions, i.e. for any $\d >0$ we have
\begin{align}  \label{S1E6} %
\int_{s(t) }^{s(t) + \d}h^3 h_{xxx}^2 \ dx  \ = \ \infty.
\end{align}%
At a first glance this would indicate that the type (b) solutions cannot have a
physical meaning and therefore must be discharged. However, the situation is
more complicated: We note that the contact angle $\ang_{\eps }$ is determined by
the classical Young's condition, and hence depends on the interfacial energies of
liquid--solid, liquid-gas and solid--gas interfaces. If $\ang_{\eps }$ diverges
logarithmically, then the energy related to liquid--solid interfaces is much
larger than the energy related to solid--gas interfaces (with a logarithmic
divergence). In the dissipation of energy formula \eqref{diss-intro} the
dissipation is finite for each $\eps >0,$ but it diverges logarithmically as
$\eps \to 0$ if $\ang_\eps \sim |\ln \eps|^{\frac 13}$. This dissipation
compensates for the increase of energy that results from the reduction of the
support of the droplet due to the reduction of the solid-liquid interface. Thus,
the divergent contributions of this term and the one due to the dissipation term
on the right--hand side cancel out in leading order as $\eps \to 0$. This
cancellation of the increase of energy due to the reduction of the support of
the droplet with the large dissipation will be described in detail in Subsection
\ref{sec-regime}. This cancellation explains why it is possible to have
physically admissible limiting solutions for which \eqref{S1E6} is divergent.

\medskip

Our formal arguments suggest that these three types of solutions are the only
possible solutions of \eqref{tfe-noslip} in the absence of fluxes across the
contact point. We give some formal arguments for the existence of these three
types of solutions for equation \eqref{tfe-noslip} and we will also give some
arguments how these solutions appear as a limit for solutions of
\eqref{tfe-slip}. A rigorous proof for the existence of type (a) solution will
be given in the upcoming paper \cite{GKV-preprint}. We also note that the
solutions of types (a), (b) and (c) are consistent with the lubrication
approximation.  This is not clear a priori in the case of the type (b) and (c)
solutions, because the angle between the liquid-gas interface and the
liquid--solid interface is large which would not be compatible with the
lubrication approximation, that requires to have nearly one-dimensional
flows. However, a more careful analysis shows that this is not the case for the
initial variables (cf. Remark \ref{rem-consist}).

\medskip

  
  % \paragraph{Plan of the paper.} The plan of this paper is the following. In
  % Section 2 we will recall the derivation of the equations \eqref{tfe-noslip},
  % \eqref{tfe-slip} by the so-called lubrication approximation. In Section
  % \ref{sec-noslip}, we discuss several classes of solutions for
  % \eqref{tfe-noslip}. In Section \ref{sec-slip}, we discuss how these solutions
  % arise as limits of \eqref{tfe-slip}.  One of the results that we prove in
  % Section 3 is that \eqref{tfe-noslip} cannot be solved in a domain with
  % changing support (in a smooth manner) if we assume that the contact angle is
  % finite and positive. In particular, we will describe the type (a) solutions
  % and (b) previously described for \eqref{tfe-noslip}. In Section 4 we will
  % describe how these solutions can be obtained as limits of solutions of
  % \eqref{tfe-slip}.

\textbf{Acknowledgement.} This project was initiated during a meeting at the
Lorentz Center in Leiden in June 2023, organized by L. Giacomelli, M. Gnann,
J. Hulshof, C. Lienstromberg and S. Sonner.  In addition to discussions with the organizers
there were many discussions with several participants such as M. Bertsch,
J. Eggers, G. Grün and J.R. King. The authors also gratefully acknowledge the
support by the Deutsche Forschungsgemeinschaft (DFG) through the collaborative
research centre ``The mathematics of emerging effects'' (CRC 1720 -539309657)
and Germany's Excellence Strategy EXC2047/1-390685813. The funders had no role
in study design, analysis, decision to publish, or preparation of the
manuscript.  HK was partially supported by the German Research Foundation (DFG)
by the project \#392124319 and under Germany's Excellence Strategy – EXC-2181/1
– 390900948.

\medskip

\textbf{Competing Interests.} The author declare to have no competing
interests.


\subsection{Discussion about results and previous literature}

Much work has been  devoted to the related family of equations
\begin{align}
  h_t +(h^n h_{xxx})_x\ = \ 0\qquad \text{in $\Ome := \{h>0\}$}  \label{tfe-n}
\end{align}%
for $n > 0$. In fact, for $n \in (0,3)$ solvability of this equation can be
expected to be the same as for \eqref{tfe-slip} since only the behaviour of the
leading term $h^n h_{xxx}$ as $h\to 0$ matters. For $n \in (0,3)$ two different
types of boundary conditions have been mostly considered: In the so--called
\textit{complete wetting regime} a zero contact angle condition is assumed; in
the partial wetting regime a non--zero contact angle is assumed. The contact
angle can be fixed and determined by Young's Law (as considered in this paper),
or variable.

  \medskip

  Existence of weak solutions of \eqref{tfe-n} in the complete wetting regime
  has been shown in
  \cite{BernisFriedman-1990,BertozziPugh-1994,BerettaBertschDalpasso-1995,Gruen-2004-1}
  for $n\in (0,3)$. Furthermore, the support of solutions with zero contact
  angle cannot shrink for $n\geq \frac{5}{2}$
  \cite[eq. (0.5)]{BerettaBertschDalpasso-1995}. For $n\geq \frac 32$, also the
  positivity set cannot shrink. Furthermore, for $n\geq 4$ the support of
  solutions remains constant. For existence of classical solutions we refer to
  e.g. \cite{GiacomelliKnuepferOtto-2008,GiacomelliKnuepfer-2010,GiacomelliGnannKnuepferOtto-2014,Gnann-2015,GnannPetrache-2018,GessGnann-2020}. Existence
  of weak solutions for stochastic thin--film equations has been considered
  e.g. in \cite{FischerGruen-2018,GessGnann-2020,DareiotisEtal-2021}. Relaxation
  to the stationary solution has been considered in
  e.g. \cite{CarrilloToscani-2002,Esselborn-2016}. Existence of weak solutions
  in the partial wetting regime has been established by Otto in \cite{Otto-1998}
  for $n=1$ and by Bertsch, Giacomelli and Karali in
  \cite{BertschGiacomelliKarali-2005} for $n\in (0,3)$. Finite speed of
  propagation is e.g.  shown in
  \cite{Bernis-1996-1,Bernis-1996-2,HulshofShishkov-1998,Gruen-2003}. For
  short--time, very initial data lead to a waiting time phenomenon, see e.g.
  \cite{DalpassoGiacomelliGruen-2001,BloweyKingLangdon-2007,Shishkov2007,Fischer-2014}. Existence
  of classical solutions with stationary support for $n = 3$ is work in
  preparation \cite{GKV-preprint}. At the contact point, these solutions have
  the behaviour $h(x,t)\approx b(t)(x-s(t))$, where $s(t)$ is the contact point.
  The long--time asymptotics of complete wetting solutions has been formally
  studied by Giacomelli, Gnann and Peschka in
  \cite{GiacomelliGnannPeschka-2023}.  Travelling wave solutions have been
  systematically investigated by Boatto, Kadanoff and Olla
  \cite{BoattoEtal-1993,KingTaranets-2013} and by King and Taranets in
  \cite{KingTaranets-2013}.

  \medskip

  The speed of the contact point for the model \eqref{tfe-slip} and its
    relation to the macroscopic or apparent contact angle has been denoted as
    Cox--Voinov Law, we refer to the works by Voinov, Tanner, Cox and Hocking
    \cite{Voinov-1977,Tanner-1979, Cox-1986, Hocking-1992}, see also
    \cite{BonnEggersEtal-2009}. In \cite{Eggers-2004}, Eggers has derived higher
    order corrections to this Law. Existence of travelling wave solutions which
    satisfy this law has been shown by Giacomelli, Gnann and Otto in
    \cite{GiacomelliGnannOtto-2016} in the zero contact angle case and by Gnann
    and Wisse \cite{GnannWisse-2022} in the case of non--zero contact
    angle. Delgadino and Mellet show that solutions of the thin--film equation
    equation in the regime of small slip converge to solutions of a quasistatic
    evolution problem \cite{DelgadinoMellet-2021}.

\medskip

The quantitative behaviour near the contact line can be best seen by the ODE 
\begin{align} \label{main-ODE} %
  ((h^3+ \eps^{3-n} h^{n})h_{\xi\xi\xi})_\xi  \ %
  = \ \pm 1\qquad \text{for $\xi \in  (0,\infty )$},
\end{align}%
which emerges e.g. from a travelling wave solution ansatz. The positive sign on
the right--hand side means a shrinking droplet, while the negative sign models
an expanding droplet. This equation for $n=3$ and $\eps =0$ has been explored
e.g. in \cite[eqs. (4.3.9),(4.3.23)]{BenderOrszag-Book}. We note that
\eqref{main-ODE} for values $n<3$ has been investigated in \cite[$r = 1$
corresponds to $m = 3$]{BoattoEtal-1993}. The equation for $n=3$ has also been
investigated in \cite{DuffyWilson-1997} in the context of Tanner's Law, relating
radius of the droplet with the time.

\medskip

Eggers and Fontelos \cite{EggersFontelos-2025} have recently considered the case of a pinned contact line
both analytically and numerically. They show that the
for $n \neq 3$, the contact angle changes like a power law, while the dynamics
become nonlocal in the critical case $n = 3$. With this result, the authors
conjecture that the contact line problem with prescribed contact angle becomes
ill--posed in the case $n = 3$.

\medskip

In addition to assuming a slip at the interface, other physical mechanisms have
been proposed as well to resolve the no--slip paradox, for reviews we refer
e.g. to \cite%
{BDQ-Book,OronDavisBankoff-1997,BonnEggersEtal-2009,SnoeijerAndreotti-2013,Shikhmurzaev-2020,DurastantiGiacomelli-2024}.
For the situation when there is evaporation of the liquid at the contact line,
we refer to \cite{RednikovColinet-2020}.  The results in that paper suggest the
existence of another class of solutions having non--trivial fluxes. Actually,
the presence of fluxes at the contact line would imply
$h\approx (x-s(t))^{\frac 34}$ near the contact point.  For computational
approaches to spreading of thin--films we refer e.g. to
\cite{BarrettBloweyGarcke-1998,MuenchWagnerWitelski-2025,Peschka-2015,Peschka-2018,PeschkaHeltai-2022}.

\medskip

\textbf{Notation.}  We will use the notation for asymptotic formulas, i.e. we
write $%
f(x) \approx g(x)$ as $x\to x_{0}\in [ -\infty ,\infty] $ to denote that
$\lim_{x\to x_{0}}\frac{%
  f(x) }{g(x) } =1$ and $f(x) \ll g(x)$ as $x\to x_{0}\in [ -\infty ,\infty] $
if $%
\lim_{x\to x_{0}}\frac{f(x) }{g(x)} =0$. We write $f(x) \sim g(x)$ if there are
constants $c,C > 0$ such that $c < \frac{f(x) }{g(x)} < C$ for $|x-x_0| \ll
1$. We will denote as $C_{x,t}^{k,m}$ the set of non-negative functions $h$ that
have $k$ derivatives in $x$ and $m$ derivatives in $t$ in the set
$\left\{ h>0\right\} $.  Notice that in this definition no assumptions are made
about the behaviour of the function or its derivatives as $(x,t)$ approaches
$\partial \left\{ h>0\right\}$. We repeatedly use the Taylor expansion
$(1-a)^{\frac 13} \approx 1 - \frac 13 a$ for $a \ll 1$


\section{Physical background \label{LubApp}}

\begin{figure}
  \centering
  \includegraphics[width=6cm]{droplet.png}
  \caption{Thickness profile $h(t,x)$ of an evolving droplet.}
  \label{fig-droplet}
\end{figure}

\subsection{Lubrication approximation} \label{ss-lubapp} %

We recall the lubrication approximation that is used to derive the equation
\eqref{tfe-slip} for $n = 2$ taking as starting point the free boundary problem
for the Stokes or Navier-Stokes equations describing the evolution of a droplet
on a solid substrate (see Fig. \ref{fig-droplet}). Since in the derivation of
the thin film equation, the inertial terms give a negligible contribution, we
will consider only the free boundary problem for the Stokes problem. For
simplicity, we consider the two-dimensional case. Since this derivation is
standard we will just indicate the main steps, see also
\cite{GiacomelliOtto-2003,GuentherProkert-2008,MatiocProkert-2012,KnuepferMasmoudi-2013,KnuepferMasmoudi-2015}. We
assume that the geometry of the fluid domain is
\begin{align}
  \UU(T) \ := \ \big \{ X \in (S_1,S_2), \ 0 < Y < H(X,T)  \big \},
\end{align}
where $H$ is the droplet height, satisfying $\spt H = [S_1(T),S_2(T)]$ also
including the case $S_2(T) = \infty$. We write
$\p_1 \UU := \OL \UU \cap \{ Y > 0 \}$ for the liquid--gas interface and
$\p_0 \UU := \OL \UU \cap \{ Y = 0 \}$ for the solid--liquid interface. The
evolution is given by the free boundary problem driven by the Stokes flow
\begin{align+}
  \mu \Delta V \ %
  &= \nabla P, %
  &\nabla \cdot V %
  &= \ 0\ \ %
  &&\text{in\ \ }\UU(t)   \label{S2E1} \\
  V_1 \ %
  &= \ \bet^n |\p_Y V_1|^{n-2} \p_Y V_1, %
  &V_{2}  %
  &= \ 0  %
  &&\text{on $\p_0 \UU(t)$},   \label{S2E3} \\
  2 \mu \SS N \ %
  &= (P - \gam_{LG} \kap) N  %
  &&&&\text{on $\p_1 \UU(t)$},  \label{S2E5} \\
  V_N \ %
  &= \ V\cdot N 
  &&&&\text{on $\p_1 \UU(t)$},  \label{S2E5b} \\
  H_X \ %
  &= \ \tan \Theta \ %
  &&&&\text{at $\p_0 \UU \cap \p_1 \UU$}, \label{S2E6}
\end{align+}
where $\SS := \frac 12 (\nabla V + \nabla^t V)$. Here, $V(X,T)$ is the fluid
velocity, $P(X,T)$ is the pressure, $\kap$ is the curvature of the interface,
$N$ its normal and $V_N$ is the normal velocity of the interface. Furthermore,
$\mu$ is the viscosity, $\gam_{LG}$ the surface tension of the liquid-gas
interface and the (generalized) Navier slip length $\bet$. On the liquid--gas
interface, we assume that the viscous stress is balanced by the surface tension
and the interface moves with the fluid. At the liquid--solid interface, we
assume a Navier slip boundary condition and a non--penetration condition. The
slip length $\bet$ is assumed to be small and noting that in the limit
$\bet \to 0$ we obtain the classical no-slip boundary condition. The system of
equations \eqref{S2E1}--\eqref{quot} is complete and well--posedness has been
shown e.g. in \cite{GuoTice-2018,LeoniTice-2023,BGKMRS-2024b}, either
for static or corresponding dynamic contact angle conditions.

\medskip

We assume the contact angle $\Theta$ to be fixed and determined by Young's Law
\begin{align} \label{quot} %
  \cos \Theta \ %
  = \ \frac{\gam_{SG}-\gam_{SL}}{\gam_{LG}} \ \in \ (0,1), %
\end{align}
which ensures local minimization of interfacial energy at the contact
point. Here, $\gam_{SL}$, $\gam_{SG}$ and $\gam_{LG}$ are the surface tensions
between solid, liquid and gas, respectively. The angle $\Theta$ is also called
microscopic contact angle which might differ from the experimentally observed
macroscopic contact angle. Solutions dissipate the total interfacial energy,
\begin{align} \label{energy-init} %
  \EE[H] \ %
  = \ \gam_{LG} \int_{\p_0 \UU} \sqrt{1+H_{X}^2} -\cos \Theta \ dX.
\end{align}
In this paper, we mostly consider the partial wetting regime where the right
hand side of \eqref{quot} is in the interval $(0,1)$ and hence
$\Theta \in (0,\pi)$.

We sketch the formal derivation of the thin film equation \eqref{tfe-slip} using
the lubrication approximation starting from \eqref {S2E1}--\eqref{S2E6}. We
first introduce a typical vertical length scale $\LAM_y$ and a typical
horizontal length scale $\LAM_x$: We choose $\LAM_y$ as the maximal droplet
thickness, assuming that this number is finite. Furthermore, we define $\LAM_x$
as a typical length scale where the change of the height from the contact point
is of the same order as $\LAM_y$.

\medskip

We make the following assumptions:
\begin{enumerate}
\item Horizontal length scales are much larger than vertical length scales. In
  terms of the non--dimensional variable this implies
  \begin{align}
  \delta \ := \ %
  \frac {\LAM_y}{\LAM_x} \ \ll  \ 1, \qquad\qquad %
  \end{align}
  We assume that a corresponding relation also holds for the derivatives of the
  profile, e.g. $|H_{XX}| \lesssim \LAM_y/\LAM_x^2$.
\item The evolution is driven by a balance of capillary and viscous forces. The
  dominance of viscous forces has been already used by the choice of Stokes flow
  for the initial model
\item We also assume that the maximal thickness of the droplet is much larger
  than the slip length, i.e. we consider the regime
  \begin{align}
  \eps \ := \ \frac{\bet}{\LAM_y}.
  \end{align}
  This assumption is, strictly speaking, not necessary for the validity of the
  lubrication approximation but is valid in the regime $\eps \ll 1$ which we
  consider in the following. We note that all parameters in the following can be
  thought as a function of $\eps$, i.e. $\d = \d_\eps$.
\end{enumerate}
We also choose  typical length scales $\LAM_p$ and $\LAM_t$ for pressure in time by
\begin{align} \label{res-2} %
  \LAM_p \ %
  := \ \frac{\gam_{LG} \LAM_y}{\LAM_x} \ %
  = \ \gam_{LG} \delta, \quad\qquad  %
  \LAM_t \ %
  := \ \frac{\gam_{LG} \LAM_y}{\LAM_x \mu}  \ %
  = \ \frac{\gam_{LG} \delta}{\mu}.
\end{align}
We introduce the rescaled (non--dimensional) variables
\begin{align} \label{res-var} %
  \big( x, y, h, t, v_1, v_2, p, \ang, \Ome \big) \  %
  := \ \Big( \frac{X}{\LAM_x}, \frac{Y}{\LAM_y}, \frac{H}{\LAM_y}, \frac{T}{\LAM_t}, \ %
   \frac{V_1 \LAM_t}{\LAM_x}, \frac{V_2 \LAM_t}{\LAM_y}, \ \frac{P}{\LAM_p}, \frac{\Theta }{\delta }, \frac{\p_0 \UU}{\LAM_x} \Big). %
\end{align}
By our choice of variables, solutions of the limit equation change in a time
scale of order $1$, see Remark \ref{rem-time}. A Taylor expansion of the stress
tensor and of the normal and tangent vectors in the new variables yields
 \begin{align}
   \SS \ %
   &\approx \ \d^2 %
   \begin{pmatrix}
     \d^2 v_{1,x} & \frac 12 v_{1,y} \vspace{0.5ex}\\
     \frac 1{2} v_{1,y} & \d^2 v_{2,y}
   \end{pmatrix}, \qquad
     N \ \approx \ \VEC{- \d h_x \\ 1}, \qquad
   N^\perp \approx \VEC{1 \\  \d h_x}.
 \end{align}
 For the boundary conditions at the upper boundary we then get
 \begin{align}
   \d^2 (N^\perp \cdot \SS N) \ %
   &= \  - \d h_x v_{1,x} + \frac 12 v_{1,y} (1 + \d^2 h_x^2) + \d h_x v_{2,y}, \\
   \d^2 (N \cdot \SS N) \ %
   &= \   \d^4 h_x^2 v_{1,x}^2 + \d h_x  v_{1,y}  + \d^2 v_{2,y}.
 \end{align}
 From the boundary condition \eqref{S2E5} and using \eqref{res-2} we hence get
 $v_{1,y} = 0$ and $p = h_{xx}$. The second momentum equation yields $p_y = 0$,
 i.e.  $p = h_{xxx}$. The first momentum equation then reads in highest order
 $v_{1,yy} = p_x = h_{xxx}$. Together with the boundary conditions
 $v_1 = \eps v_{1,y}$ for $y = 0$ and $v_{1,y} = 0$ for $y = h$ we obtain the
 velocity profile
 \begin{align}
   \mu v_1 \ = \ h_{xxx} \Big (\frac 12 y^2 - h y - \eps h \Big) \qquad \text{for $y \in (0,h)$. }
 \end{align}
 The equation for conservation of mass can be written in the form
 $\p_t h + v_1 h_x = v_2$. Integrating in vertival direction, we then obtain the
 thin--film equation
 \begin{align}
   h_t \ %
   = \ - \Big (\int_0^h  v_1 \ dy \Big)_x \ %
   = \ - \Big (\big(\frac 16 h^3 + \eps h^2 \big)  h_{xxx} \Big)_x. \ %   
 \end{align}
 The factor $\frac 16$ can be removed by a further simple rescaling in $h$ and time.

 \medskip

 It might seem inconsistent that we allow for large contact angles in the
 lubrication approximation. However, this contradiction is resolved by the
 difference between physical and rescaled variables as indicated in Section
 \ref{LubApp}.
 \begin{remark}[Applicability and rescaled contact angle] \label{rem-consist} %
   \text{} %
   For the lubrication approximation to be a good approximation, the droplet
   profile has to be flat and in particular the contact angle has to be small in
   the initial variables.  Due to the anisotropic rescaling, the droplet is,
   however, not necessarily flat in the rescaled variables and in particular the
   rescaled contact angle $\ang$ is not necessarily small. In fact, we consider
   the limit $\eps \to 0$ where the angle $\ang_\eps$ is a second free parameter
   which can take different values in the limit. The smallness of $\delta$ then
   ensures the flatness of the liquid droplet.
\end{remark}
Since the contact angle is determined by Young's Law we can also phrase the
regimes (a),(b) and (c) from the introduction in terms of the interfacial
energies.
\begin{remark}[Lubrication approximation and interfacial energies]
  For the validity of the lubrication approximation we need
  $\frac 12 \Theta^2 \approx 1 - \cos \Theta$. In view of Young's Law
  \eqref{quot} this is the same as
  \begin{align}
      \frac 1{\gam_{LG}} (\gam_{LG} + \gam_{SL} - \gam_{SG})  \ %
      \approx  \ \frac {\Theta^2}{2} \ %
      = \ \frac 12 \ang^2 \d^2 \ %
      \to \ 0^+ \ \qquad %
      \text{as\ $\eps \to 0$}  \label{S4E2}
  \end{align}
  Type (b) solutions are then e.g. obtained in the regime when
  \begin{align}
      \frac 1{\gam_{LG}} (\gam_{LG} + \gam_{SL} - \gam_{SG})  \ %
    \approx \ \gam \d^2 |\ln \eps|^{\frac{2}{3}} \qquad\qquad %
   \text{as\ $\eps \to 0$}  \label{S4E3}    
  \end{align}
  In particular, the interfacial energies $\gam_{SG}$, $\gam_{SL}$, $\gam_{LG}$
  need to be adjusted in a very careful way to obtain the dynamics of the
  interface by means of \eqref{tfe-slip}. In particular, type (b) solutions
  appear only if condition \eqref{S4E3} holds.
\end{remark}

\subsection{Dissipation of energy  \label{sec-regime}}

Solutions of the initial Stokes system dissipate the total interfacial energy
\eqref{energy-init}. With a small angle expansion of the energy
\eqref{energy-init} and in terms of the rescaled variables \eqref{res-var} we get
\begin{align}
  \frac{\EE[H]}{\gam_{LG} \LAM_x \d^2} \  %
  \lupref{energy-init} = \ \frac 1{\LAM_x\d^2} \int_{\p_0 \UU} \sqrt{%
  1+H_{X}^{2}} -\cos \Theta \ dX  \ %
  \approx  \
  \FF[h],
\end{align}
where the thin--film energy is given by
  \begin{align}
    \FF[h] \ := \ \frac 12  \int_{\Ome} h_x^2+ \ang^2 \ dx.
  \end{align}
  The thin film equation \eqref{tfe-slip} satisfies an energy--dissipation
  identity with respect to quantity on the right--hand side above even for weak
  solutions.  A definition for weak solutions in the partial wetting regime has
  been given by Bertsch, Giacomelli and Shishkov in
  \cite[Def. 3.1]{GiacomelliShishkov-2005}. Here, we recall this identity for
  classical solutions with sufficient regularity. In fact, let $n\in (0,3)$,
  $\eps > 0$ and let $h$ be a sufficiently regular solution of
  \eqref{tfe-slip}--\eqref{tfe-bc}.  Then
  \begin{align}
      \ddt \FF[h] \ = \ - \int_{\Ome} (h^3+ \eps^{3-n} h^{n}) h_{xxx}^2 \ dx. \label{T5E4}
  \end{align}
  To see these we may assume without loss of generality we assume $\Ome=(s_{-},s_{+})$. Using
  \eqref{tfe-bc} we get
\begin{align}
  \frac d{dt} \frac 12 \int_{\Ome} h_x^2 + \ang_\eps^2  \ dx   \ %
  = \ \int_{\Ome}h_{tx}h_x=-\int_{\Ome}((h^3+ \eps^{3-n} h^{n}) h_{xxx})_{xx}h_x \ dx.
\end{align}
Integrating by parts two more times yields \eqref{T5E4}. The boundary terms for
the integration by parts for the last identity vanish by assumption on the
regularity of the solution. We now discuss the dissipation of energy for type
(a) solutions.
\begin{remark}[Dissipation of energy for type (a) solutions] %
  Suppose that $\Ome (t)$ is a bounded interval. In this case, formal analysis
  Tanner (see Section \ref{ss-TannerCV} below) indicates that
  \begin{itemize}
  \item in the first stage where $t = \OO(1)$, $|\Ome (t) |$ remains nearly
    constant, and most of the dissipation takes place through the term
    $\int_{\Ome (t)}h^3h_{xxx}^2dx$ via reduction of the energy
    $\int_{\Ome (t)}h_x^2dx$.
      \item During the second stage where $t = \OO(|\ln \eps|$), $h_{xxx}$ is
        approximately zero in the bulk of the droplet, and most of the dissipation
        takes place at $\OO(\eps)$--distance from $\p \Ome (t)$. This dissipation rate
        is of order $|\ln \eps|^{-1}$ and is associated to $\OO(1)$ changes of
        $|\Ome (t)|$ in times of order $|\ln \eps|$. This can be seen with the change
        of variables $\eps y = x-s$, $h= \eps H$ (see \eqref{S7E4}).
      \end{itemize}
    \end{remark}
    Type (b) solutions are obtained in the regime
    $\ang_\eps \approx \gam |\ln \eps|^{\frac 13} \ll 1$. For type (b) solutions
    the relaxes energy by reducing the size of the solid--liquid interface. The
    energy gained in this way is infinitely (logarithmically) large and it is
    compensated with the infinite (logarithmic) dissipation in the fluid.
  \begin{lemma}[Energy dissipation for type (b) solutions]
    Let $h$ be a solution of \eqref{tfe-noslip} with support $(s,\infty)$ for
    some function $s = s(t)$. We assume that $h$ is a type (b) in the sense that
    for $\xi :=x- s$ we have
    \begin{align}
      h(\xi,t) \ %
      \approx \ (3 \dot s)^{\frac 13} \xi |\ln \xi|^{\frac 13} \qquad
      \text{for $\xi \ll 1$.}
    \end{align}
    We also assume that the time derivative $h_t$ and up to three spatial
    derivatives of $h$ have the asymptotic behaviour, obtained by applying these
    derivatives on the right--hand side above.  Furthermore, the solution and its
    derivatives decay sufficiently fast for $\xi\gg 1$. Then the limit below
    exists and is finite and
    \begin{align}
      \frac d{dt} \frac 12   \int_{s}^\infty h_x^2\ dx    \ %
      = \ \lim_{\d \to 0} \Big( \frac{\dot s}2h_x^2(\d )-\int_{s + \d}^{\infty}h^3h_{xxx}^2 \ dx \Big).
    \end{align}
\end{lemma}
\begin{proof}
  By assumption, for $\xi \ll 1$ we have
  $h_{\xi\xi\xi} \approx \frac 13 (3 \dot s)^{\frac 13} \xi^{-2} |\ln \xi|^{-\frac 23}$ and
  \begin{align}  %
    &h_\xi h_t  \ %
    = \ \OO(\xi |\ln \xi |^{\frac23}) \ = \ o(1), \label{asy1} \\ \qquad %
    &h^3h_{\xi\xi}h_{\xi\xi\xi} \ %
    = \ \OO(|\ln \xi |^{-\frac 13}) \ = \ o(1). \label{asy3} %
  \end{align}
  where $o(1) \to 0$ as $\xi \to 0$.  For $\xi > 0$ we have
  $h_t = h_\xi\dot s - (h^3 h_{\xi\xi\xi})_\xi$. Using this identity, for
  $\d > 0$ we hence obtain
  \begin{align}
    \ddt \frac 12 \int_{\d}^{\infty } h_\xi ^2\ dx  \ %
    &\stackrel{p.i.}= \ -\int_{\d}^{\infty }h_{\xi\xi}h_t + \left[ h_{\xi
      }h_t \right]_{\d}^{\infty } \notag \\
    &\upref{asy1}= \ -\dot s\int_{\d}^{\infty } h_{\xi\xi} h_\xi  \ d\xi %
      + \int_{\d}^{\infty }h_{\xi\xi}(h^3h_{\xi\xi \xi})_\xi +  o(1) \notag \\
    &= \ \frac{\dot s}2h_\xi ^2(\d )-\int_{\d}^{\infty} h^3h_{\xi\xi \xi}^2 \ d\xi %
      -\left[ h^3h_{\xi\xi}h_{\xi\xi \xi}  \right]_{\d}^{\infty }+ o(1) \notag \\
    &\upref{asy3}= \ \frac{\dot s}2h_\xi ^2(\d) - \int_{\d}^{\infty}h^3h_{\xi\xi \xi}^2 \ d\xi  + o(1),  \notag
  \end{align}%
  where $o(1) \to 0$ as $\d \to 0$. For $\d \to 0$ we obtain the asserted
  identity. Furthermore, we note that we have  the integral identity
  \begin{align}
    \int_\d^{\frac 12} \frac 1{\xi |\ln \xi|^{\frac 13}} \ %
    \approx \ \frac 32 |\ln \d|^{\frac 23}  \qquad\qquad \text{for $\d \ll 1$}
  \end{align}
  In particular, for  $\xi \to 0$ a straightforward calculation yields
  \begin{align}
    \frac{\dot s}2h_\xi ^2(\d) - \int_{\d}^{\infty}h^3h_{\xi\xi \xi}^2 \ d\xi \ %
    \approx \ \frac 16  (3\dot s)^{\frac 53} |\ln \d|^{\frac 23}  -  \frac 16 (3 \dot s)^{\frac 53} |\ln \d|^{\frac 23} + \OO(1).
  \end{align}
  i.e. the cancellation of the leading order terms shows that the limit in
  leading order exists and is finite.
\end{proof}
Generally, the interfacial energies between liquid and solid might depend on the
horizontal variable. The contact angle then also depends on the position.  In
this case, for the rescaled model and for type (b) solutions we would get a
contact angle of the form
$\ang_{\eps}(x) = \psi (x) |\ln \eps|^{\frac 13}\to \infty$ for some function
$\psi = \psi (x)$,$\,$i.e. $\psi (x)$ would modify the angle locally. The energy
then has the form
\begin{align}
  E(h) \ %
  = \ \int_{s(t)}^{\infty }h_x^2 \ dx +|\ln \eps|^{\frac23} \int_{s(t)}^{\infty }\psi (x)\ dx.
\end{align}


\section{Solutions in the case of no slip} \label{sec-noslip} %

We now discuss some particular problem classes for equation
\eqref{tfe-noslip}. We will consider two types of problems, namely, solutions
with shrinking support and solutions with fixed support.

\subsection{Sufficient conditions for non--spreading of solutions}

We consider sufficient conditions for the non--spreading of solutions. Since we
do not consider the asymptotic behaviour $\eps \to 0$ here, we can equivalently
discuss the issue for the equation \eqref{tfe-n} keeping only the dominant term
near the boundary. It is not possible to have smooth solutions with a finite
contact angle if the interface moves. In this section, we consider local
solutions in the domain
\begin{align}
  D \  =  \ \left\{ (x,t) :s(t) <x<s(t) + \d ,\ 0<t<T\right\} 
\end{align}
for some $\d > 0$ and a single contact point $x = s(t)$, i.e. we consider
\begin{align} \label{sol-local} %
  \begin{aligned}
    &h_t + (h^n h_{xxx})_x \ = \ 0 \qquad &&\text{in $D$}, \\
    &h \ = \ 0 \qquad &&\text{for $x = s(t)$}.
  \end{aligned}
\end{align}
The no--slip paradox is the conjecture that the no--slip condition at the
liquid--solid interface prevents movement of the contact line since this would
lead to infinite dissipation of energy. In the perspective of the thin--film
equation we can phrase it in the following way:
\begin{conjecture}[Finite dissipation implies fixed support] %
  Let $n \geq 3$, $\d >0,\ T>0$. Let $h\in C^0(\OL D)$ and $s\in C^0((0,T))$ be
  a classical and positive in $D$ solution of \eqref{sol-local} with %
  \begin{align} \label{S1E8} %
    \iint_{D} h^n  h_{xxx}^{2} \ dxdt \ < \ \infty.  
  \end{align}
  Then $s(t) = s_0$ is constant in $(0,T)$.
\end{conjecture}
We have assumed continuity of the height profile. Alternatively, one could also
assume finiteness of the thin--film energy.  The next theorem was discussed
during a workshop in the Lorentz Center \cite{Lorentz-Center}. In particular,
J. Eggers asked if under general smoothness conditions one can prove that the
support of the solution does not expand. For the convenience of the reader we
present a short and simple proof: 
\begin{theorem}[Regularity implies fixed support] \label{thm-nomove} %
  Let $n \geq 3$, $\d >0,\ T>0$ and let $h\in C_{x,t}^{4,1}(\OL D)$ with
  $s\in C^{1}((0,T))$ be a solution of \eqref{sol-local} such that
  \eqref{no-massflux} holds. Furthermore, suppose that for some positive
  $\gam\in C^{1}((0,T))$ and $R_1, R_2 \in C_{x,t}^{0,1}(D)$ we have
  \begin{align} \label{exp-regder}
    h(x,t) \ %
    &= \ \gam (x-s(t) ) + R_{1}(x,t) \qquad \text{where $\DS \lim_{x\to s(t)^+}\frac{R_{1}(x,t) }{x-s(t)} \ = \ 0$}, \\
    h_t (x,t) \
    &=-%
    \gam \dot s(t) +R_{2}(x,t)\ \qquad\qquad %
    \text{where  $\lim_{x\to s(t)^+}R_{2}(x,t) \ =\ 0$.} \label{exp-timeder} %
  \end{align}
  Then $s(t) = s_0$ is constant in $(0,T)$.
\end{theorem}
\begin{proof}
  The expansion \eqref{exp-timeder}  implies that for any $t \in (0,T)$ we have     \text{as $x\to s(t)^+$}
  \begin{align}
    (h^n h_{xxx})_x \ %
    = \ - h_t \ 
    \stackrel{\eqref{exp-timeder}}= \ \gam \dot s(t) - R_{2}(x,t).  % 
  \end{align}
  Integrating once from $s(t)$ to $x$ and using the assumptions  we obtain
  \begin{align}
    h_{xxx}(x,t) \ %
    &\stackrel{\eqref{no-massflux}}= \  \frac 1{h^n} \Big(\gam \dot s(t) (x-s(t)) + \int_{s(t)}^x R_2(\t x,t) \ d \t x \Big) \stackrel{\eqref{exp-regder}}\approx \  \frac{\dot s(t)}{\gam^{n-1} (x-s(t))^{n-1}}. \qquad\qquad %
  \end{align}
  Integrating two times and if $n \geq 3$ it follows that
  $|h_x(x,t)| \to \infty$ as $x\to s(t)^+$ if $\dot s(t) \neq
  0$. This implies  $\dot s = 0$.
\end{proof}
We note that the above proof does not directly use finite dissipation of
energy.

\subsection{Type (a) solutions -- fixed  support}

In this section, we justify why solutions of \eqref{tfe-noslip} with fixed
support $\supp h = [0,\infty)$ can be expected to exist. We consider the
problem%
\begin{align}  h_t +(h^3h_{xxx})_x\ %
  &= \ 0\ \ \qquad\qquad\text{for }x\in (0,\infty ), \  t>0,  \label{S4E5} \\
  h(0,t) \ %
  &= \ 0\ \ \qquad\qquad\text{for }t\geq 0  \label{S4E6}
\end{align}
with initial data $h_0$. One can construct solutions to \eqref{S4E5}--\eqref{S4E6} with linear behaviour
near the contact point. More precisely, we assume that
\begin{align}   \label{S4E9} %
  h(x,t) \ %
  = \ \gam x + \phi (x,t),
\end{align}
where $\gam$ is a smooth positive function which represents the leading order
behaviour near the contact point. A rigorous construction for solutions of
\eqref{S4E5}--\eqref{S4E6} is work in progress \cite{GKV-preprint}. Here, we
just give formal arguments to show the feasibility of such solutions.

\medskip

Plugging \eqref{S4E9} into \eqref{S4E5} and keeping the leading order terms we
obtain%
\begin{align}
  \phi_t + \gam^3(x^3\phi_{xxx})_x + \dot{\gam} x \ %
  \approx \ 0.
\end{align}
Using that $\gam >0$ we can change the time scale as
$\tau = \int_{0}^{t}(\gam(s) )^3ds$. Then%
\begin{align}
  \phi_{\tau }+(x^3\phi_{xxx})_x + \frac{\dot{\gam} x}{\gam^3} \ = \ 0 \qquad\qquad %
  \text{for $x>0$,} \label{S5E1}
\end{align}%
where $\dot \gam = \ddt \gam$. The leading order behaviour for $x \ll 1$ is
obtained by%
\begin{align}
  \phi_{\tau} + (x^3\phi_{xxx})_x \ = \ 0 \ \qquad %
  \text{for $x>0$.}
\label{S5E3}
\end{align}
Equation \eqref{S5E3} is a fourth order parabolic equation degenerated at
$x=0$. Let $\tilde{\phi}$ denote the Laplace transform of $\phi$ in time with
Laplace variable $z \in \C$. Then \eqref{S5E3} turns into the ODE %
\begin{align} \label{this-ODE} %
  z \tilde{\phi}+(x^3\tilde{\phi}_{xxx})_x \ %
  = \ 0\qquad\qquad %
  \text{for $x>0$, for all $z \in \C$}.
\end{align}
The possible asymptotic behaviours as $x\to 0^+$ for solutions of
\eqref{this-ODE} are $a_1 \ln x$, $a_2$, $a_3 x$, $a_4 x^{2}$ for suitable
functions $a_i = a_i(z)$. Since we consider solutions for which the contact
point stays at $x=0$ we must impose $a_1(z) =a_{2}(z) =0$.  Therefore, the
solution of \eqref{S5E3} behaves to the leading order as $%
\phi \approx A_1(\tau ) x$ as $%
x\to 0^+$. On the other hand, the problem in which we are interested is
\eqref{S5E1}. We can then choose $\gam$ such that $\phi$ in \eqref{S4E9}
satisfies $\phi =O(x^{2})$ as $x\to 0^+$. Then $h \approx \gam x$ as
$x \to 0^+$.

\subsection{Type (b) solutions -- prescribed shrinking support \label{IncreasingS}}

We now discuss a different class of solutions of \eqref{tfe-noslip}, supported
in the half-line $(s(t) ,\infty)$ for $t \geq 0$. The remarkable feature is that
the position of the contact point $s(t)$ can be chosen arbitrarily, as long as
it is sufficiently smooth and strictly increasing. We assume
$\supp h = [s(t),\infty)$ i.e. there is only a single contact point. In terms of
$\xi = x - s(t)$ \eqref{tfe-noslip} takes the form
  \begin{align} \label{tfe-noslip-xi} %
    h_t  - \dot s h_\xi +(h^3 h_{\xi\xi\xi})_\xi \ &= \ 0  \qquad 
    \text{for $\xi > 0$}, \\
    h(0,t) \ &= \ 0
  \end{align}
  with initial data $h(\cdot,0) = h_0$. We write by some abuse of notation $h$
  as a function of $\xi$ and $t$. We look for solutions of \eqref{tfe-noslip-xi}
  with the asymptotic behaviour
\begin{align}   \label{ass-type-b} %
  h(\xi,t) \ \approx \ \gam \xi  (- \ln \xi)^{\frac 13} \qquad\qquad %
  \text{as } \xi \ll 1.
\end{align}
Plugging the asymptotics \eqref{ass-type-b} into \eqref{tfe-noslip-xi} we obtain
the relation $3 \dot s \ = \ \gam^3$ i.e. we get the
asymptotics \label{def-typeb} of the type (b) solutions. The non--negativity of
the solution together with \eqref{ass-type-b} implies the consistency condition
$\dot s >0$. Solutions with the asymptotics \eqref{ass-type-b} have been
investigated for travelling solutions e.g. in
\cite{BoattoEtal-1993,BowenKing-2001}. We remark that these solutions do not
have a finite contact angle.

\medskip

Linearizing formally around the asymptotic behaviour \eqref{ass-type-b} we
obtain a linear equation near the contact point which is solvable in the class
of functions with linear behaviour including a logarithmic correction. We write
\begin{align}   \label{H-bar-def} %
  h(\xi,t) \ %
  =: \ \OL{h} ( \xi,t) +u( \xi,t), \qquad\qquad %
  \text{where} \quad %
  \OL{h}(\xi,t) \ %
  := \ \gamma \xi ( - \ln  \xi) ^{\frac 13}%
\end{align}%
assuming that $u(\xi,\cdot) \ll \OL{h}(\xi,\cdot)$ for $\xi \ll 1$ with
similar estimates for spatial derivatives. In leading order we obtain the
approximated equation as $\xi\to 0^+$
\begin{align}
  \OL{h}_t +u_t + \big[ ( \OL{h}^{3}\OL{h}_{\xi\xi\xi})_\xi -\dot s \OL{h}_{\xi}\big] +( \OL{h}^{3}u_{\xi\xi\xi}+3\OL{h}^{2}%
  \OL{h}_{\xi\xi\xi}u-\dot s   u)_\xi \ %
  = \ 0.  \label{u-nonhomog}
\end{align}
The leading order behavior of $u$ as $\xi\to 0^+$ can be computed arguing
similarly as in the case of fixed, prescribed support (Subsection 3.2). We need
to examine the contribution of the homogeneous terms, namely%
\begin{align}
u_t + \big( \OL{h}^{3}u_{\xi\xi\xi}+3\OL{h}^{2}\OL{h}_{\xi\xi\xi}u-\dot s  u\big)_\xi \ %
= \ 0.  \label{u-homog}
\end{align}
A major difference with the case of constant support is the fact that the term
$\dot s u$ yields a contribution comparable to the term
$\OL{h}^{3}u_{\xi\xi\xi}$. The leading order terms for $\xi\ll 1$ in
\eqref{u-homog} are%
\begin{align}
  u_t + \gamma \big( \xi^{3} (- \ln \xi) u_{\xi\xi\xi}+ \frac{2}{3}u\big)_\xi \ %
  = \ 0 \qquad %
  \text{for $\xi \ll 1$}.
\label{u-homog-app}
\end{align}
The four possible asymptotic behaviors for $\xi \ll 1$ for solutions of
\eqref{u-homog-app} are %
$a_1 (- \ln \xi) ^{\frac 13}$, $a_2$, $a_3 \xi (- \ln \xi)^{-\frac{2}{3}}$ and
$a_4 \xi^{2}(- \ln \xi)^{\frac 13}$ for suitable functions $a_i = a_i(t)$. Since
\eqref{u-nonhomog} is a fourth order equation, we can expect to determine
uniquely the solutions prescribing two boundary conditions at $\xi = 0$. With
our assumptions we get $a_1 = a_2 = 0$. Therefore, the first contribution due to
the homogeneous part of \eqref{u-nonhomog} has the behaviour
$a_3 \xi (- \ln \xi)^{-\frac{2}{3}}$ which satisfies the consistency condition
$u(\xi,\cdot) \ll \OL{h}(\xi,\cdot)$ for $\xi \ll 1$. We need to estimate also
the contribution of the non-homogeneous terms in \eqref{u-nonhomog} for
$\xi \ll 1$. The term $%
[( \OL{h}^{3}\OL{h}_{\xi\xi\xi})_\xi -\dot s (t) \OL{h}_{\xi}]$ yields a term of
form $a(t) \xi(- \ln \xi)^{-\frac{2}{3}}\ln (- \ln \xi)) $, while the term
$\OL{h}_t $ yields a much smaller contribution, which can be estimated as
$\OO(\xi^{2}(- \ln \xi) ^{\beta }) $ for $\xi \ll 1$.

\medskip

We remark that the solvability of \eqref{tfe-noslip} in a domain
$\{ (x,t) :t>0,\ x>s(t)\} $ with the behaviour \eqref{ass-type-b} can be thought
of as a boundary value problem, with unusual boundary conditions that are due to
the degeneracy of \eqref{tfe-noslip} near the contact line $x=s(t)$. Notice that
the function $s(t)$ is arbitrary. Therefore, the support of $h$ is basically
arbitrary, as long as it is decreasing in time. These solutions are the
solutions that have been termed as type (b) solutions in the
introduction. Solvability of this problem will be addressed in future work.

\section{Solutions in the case of  weak slippage} \label{sec-slip} %

We consider the equation \eqref{tfe-slip} with weak slippage. In the first
subsection we discuss well--posedness for the equation. We then discuss how the
solutions of \eqref{tfe-slip} converge to solutions of type (a),(b),(c) of the
limit equation \eqref{tfe-noslip}. For simplicity we assume
$\supp h = [s(t),\infty)$ and we use the variables $\xi = x - s(t)$. Then
\eqref{tfe-slip} takes the form
  \begin{align} \label{tfe-slip-xi} %
    h_t  - \dot s h_\xi +(( h^3 + \eps^{3-n} h^n) h_{\xi\xi\xi})_\xi \ &= \ 0  \qquad 
    \text{for $\xi > 0$}, \\
    h(0,t) \ = \ 0, \quad h_\xi(0,t)  \ &= \ \ang
  \end{align}
  with initial data $h(\cdot,0) = h_0$. By an abuse of notation we write 
  $h$ as a function of $\xi$ and $t$.


\subsection{Well--posedness and Boundary conditions}

We consider well--posedness and boundary conditions for the problem
\eqref{tfe-slip-xi}. Since we do not consider the asymptotic behaviour
$\eps \to 0$ here, we can equivalently discuss the issue for the equation
\eqref{tfe-n}

\paragraph{Solutions with prescribed contact line $s(t)$.} one can construct some solutions with prescribed position of the contact
points. For simplicity we consider the case when $\spt h = (s(t,\infty))$. The
difference is that in the case $m<3$ one can prescribe $s(t)$ as any arbitrary
smooth function, not necessarily increasing. The behaviour of the solution
$h(x,t)$ near the contact point $x=s(t)$ with $\xi := x - s(t)$ is given by
\begin{align}
  h(x,t) \ = \ \gam \xi  + \phi (\xi,t), \qquad %
\end{align}%
for some smooth function $\gam >0$ and some function $\phi$ with
\begin{align} \label{ass-phi} %
  \phi(\xi,t) = o(\xi ), \ \  %
  \phi_\xi(\xi,t) = o(1), \ \  %
  \phi_t(\xi,t) = o(\xi) \qquad %
  \text{as $\xi\to 0^+$}. 
\end{align}
We examine the behaviour of the linearized problem satisfied by $\phi $ near
the contact point. Keeping the leading order terms we obtain%
\begin{align}
  -\gam \dot s  + \phi_t -\dot s  \phi_\xi + \gam^n (\xi^n \phi_{\xi\xi\xi})_\xi  \ %
  =  \ 0 \qquad\qquad %
  \text{for $\xi >0, t>0$},
\end{align}%
with  $\phi (0,t) =0$. By assumption \eqref{no-massflux} we
have $\xi^n \phi_{\xi\xi \xi}\to 0$ as $\xi\to 0^+$. Integrating the equation in
$\xi$ and using \eqref{ass-phi} we obtain the asymptotic behaviour
\begin{align}   \label{S5E6} %
  \phi(\xi,t) \ \approx \
  \left \{
  \begin{array}{ll}
      \DS \frac{\gam^{1-n}\dot s }{(4-n)(3-n)(2-n) }\xi^{4-n}\ \ \ \qquad %
    &\text{if $2 < n < 3$,} \vspace{0.5ex}\\
    \DS \tfrac 12 \gam^{-1}\dot s  \xi^2 \ln \xi\ \ \ \qquad %
    &\text{if $n = 2$,} \vspace{0.5ex} \\
    \DS \tfrac 12 \phi_{\xi\xi}(0) \xi^2.
    &\text{if $0 < n < 2$}.
  \end{array}
  \right.
\end{align}%
Notice that there are no constraints about the sign of $\dot s$. This is
different to the case of solutions for \eqref{tfe-noslip} where the position of
the contact point can be prescribed arbitrarily as long as the support is
shrinking (cf.  Section \ref{IncreasingS}).

\paragraph{Solutions with prescribed contact line angle.} One can construct the
solution of a different type of boundary problem for \eqref{tfe-n} with
$n<3$. More precisely, we can prescribe the value of the contact angle
$\gam \geq 0$. To impose this quantity is a very stringent condition and this
determines uniquely the value of $s(t)$. Existence of such solutions has been
proved in the complete wetting as well as partial wetting case, see
e.g. \cite{BertozziPugh-1994,BerettaBertschDalpasso-1995,BertschDalpassoGarckeGruen-1998,DalpassoGarckeGruen-1998,Gruen-2004-1,GiacomelliKnuepferOtto-2008,GnannWisse-2022}. It
seems possible to use an approach similar to the one in those papers to solve
the problem with prescribed, smooth $\gam$ by reformulating the angle condition
as an integral equation for $s(t)$.
\begin{remark}[Boundary conditions] %
  As discussed before, for our boundary value problem, we can either prescribe
  two boundary conditions for $h$ or we can prescribe the contact point velocity
  and one boundary condition. We show how a similar issue arises from the
  classical one--dimensional Stefan problem
    \begin{align} \label{stefan}
      &u_t  \ = u_{xx}\ \ 
      &&x>s(t),  \ t \geq 0, \\
      &u  \ = \  1, \quad  \dot s \ = \ - u_x, 
      &&x = s(t),  t \geq 0
    \end{align}
    with initial conditions $u = u_0$. In terms of $U = \int_x^\infty u$ for
    $x > s(t)$ we obtain
    \begin{align} \label{stefan-2} 
      &U_t  \ = U_{xx}\ \ 
      &&x>s(t),  \ t \geq 0, \\
      &U  \ = \  \int_{s_0}^\infty u_0 \ dx, \quad  U_x \ = \ - 1,  
      &&x = s(t),  t \geq 0. \label{U-bc}
    \end{align}
    Indeed, for the first boundary condition in \eqref{U-bc} we calculate
    \begin{align}
      \ddt U(s(t), t) \ 
      &= \ U_x(s(t),t) \dot s  + \p_t U (s(t),t) \ 
        = \ - \dot s  + \int_{s(t)}^\infty u_{xx} (\xi,t) \ d\xi \\ 
      &= \ - \dot s  - u_{x} (\dot s,t) \ = \ 0. 
    \end{align}
    This illustrates that the Stefan problem \eqref{stefan} -- one boundary
    condition for $u$ and the speed $\dot s$ is prescribed --- is equivalent to the
boundary value problem \eqref{stefan-2} where two boundary conditions on the
function are assumed at the free boundary while the speed of the contact line is
only implicitly given.
\end{remark}

\subsection{Approximating type (a) solutions -- Cox-Voinov Law} \label{ss-TannerCV}

In this subsection we discuss the type (a) solutions and their relation to the
Cox-Voinov Law which connects the speed of the contact point with the
  apparent contact angle.  We will see that for $\eps \to 0$, type (a)
solutions of \eqref{tfe-noslip} arise as limit of solutions of
\eqref{tfe-slip-xi} where the constant angle $\ang_\eps$ does not depend on
$\eps$. Both in the case of complete or partial wetting, the Cox--Voinov Law can
be reformulated as a law for the speed of the contact point. Our main goal is to
illustrate the fact that those computations yield a limit law for the speed of
the contact lines associated to the equation \eqref{tfe-noslip} obtained as a
limit of \eqref{tfe-slip}. In this regime the contact line moves very slowly
with speed of order $|\ln \eps|^{-1}$.

\medskip

We consider solutions of \eqref{tfe-slip-xi} with a single contact point.  Far
from the contact point, the solutions approximately solve \eqref{tfe-noslip} and satisfy
\begin{align} \label{S5E8} %
  h(\xi,t) \ %
  \lupref{S1E4}\approx \ \gam \xi, \qquad\qquad \text{for $\eps \ll \xi \ll 1$},
\end{align}%
where $\gam >0$ is an $\OO(1)$ function that depends on
initial data. The asymptotic formula \eqref{S5E8} breaks down near the contact line when
$h^3 \sim \eps^{3-n}h^{n}$, i.e.  $\xi \sim \eps$. To investigate the solutions in the inner
region we use the new variables%
\begin{align} \label{inner-var-a} %
  h \ = \ \eps H,\ \quad \ \xi \ = \ \eps y,\ \quad  t \ = \ \eps \tau  
\end{align}
In terms of these new variables \eqref{tfe-slip-xi} takes the form
\begin{align}
  H_{\tau }-\dot s H_y+ \left(\left(H^3+H^{n}\right)
  H_{yyy}\right)_y \ %
  = \ 0 \qquad\qquad
  \text{for $y>0$}. 
\end{align}
with $H_y(0,\cdot) = \ang$.  We expect the solutions to become stationary if
$\tau $ is large. With the assumption \eqref{no-massflux} we obtain the
approximative equation
\begin{align}   \label{S5E9} %
  -\dot s + (H^2+H^{n-1}) H_{yyy} \ = \ 0 \qquad\qquad
  \text{for $y>0$}.
\end{align}
We need to match solutions of \eqref{S5E9} with the asymptotic behaviour
\eqref{S5E8} for $\eps \ll \xi\ll 1$ where $\dot s$ is one of the quantities
that must be determined. The equation \eqref{S5E9} must be solved with contact
angle $\ang = 0$ in the complete wetting case and contact angle $\ang > 0$ in
the partial wetting case. This analysis to the Cox--Voinox Law in the case of
complete wetting as well partial wetting. We use the self--consistent assumption
$|\dot s| \ll 1$ for $\eps \ll 1$.


\paragraph{Partial wetting case.}

We now consider solutions of \eqref{S5E9} with nonzero contact angle $\ang > 0$
assuming that the slope is almost constant for distances to the contact point of
order $\eps$. We assume $n \in (1,3)$ to be fixed and consider the asymptotics $\eps \to 0$.  We make the ansatz
$H = H_0+H_1 + \ldots$ with $|H_1| \ll H_0$ for $y \to 0^+$. We obtain to
leading order%
\begin{align}
  H  \ %
  \approx \ H_0 \ %
  = \ \ang y \qquad\qquad \text{as $y\to 0^+$}.
\end{align}
Furthermore, $H_1$ solves to leading order
\begin{align}
  H_{1,yyy} \ %
  = \ \frac{\dot s}{H_0^2+H_0^{n-1}} \ %
  = \ \frac{\dot s}{\ang^2 y^2  + \ang^{n-1} y^{n-1}} \qquad\qquad %
  \text{for $y>0$.} 
\end{align}
With the assumption $|H_1| \ll y$ for $y \ll 1$ and $H_1 \ll y^2$ for $y \gg 1$
this yields%
\begin{align} \label{H1-form} %
  H_1 \ %
  = \ |\dot s| \int_0^y \int_0^{y_1} \int_{y_2}^{\infty }\frac{1}{\ang^2 y_3^2 +
  \ang^{n-1} y_3^{n-1}} \ dy_3 dy_2 dy_1.
\end{align}
Integrating \eqref{H1-form} we get the asymptotic behaviour 
\begin{align}
  H_1 \ %
  \approx \ \frac{|\dot s|}{\ang^2}y\ln y \qquad\qquad %
  \text{as $y\to \infty$.}
\end{align}
Using the approximation $H = H_0 + H_1$ and in terms of the variables $h$,
$\xi$, cf. \eqref{inner-var-a} and as in \eqref{S7E1} we then obtain%
\begin{align} \label{Tan-outer-match} %
  h \ = \ \eps H \ %
  \approx \ \ang y+ \frac{|\dot s|}{\ang^2}y\ln y \ %
  \approx \ \ang\xi + \frac{|\dot s|}{\ang^2} \big(\ln \frac 1{\eps}\big) \xi  \qquad\qquad %
  \text{for $\eps \ll \xi \ll 1$}
\end{align}
Matching \eqref{Tan-outer-match} with \eqref{S5E8} we obtain the Cox--Voinov Law
\cite{Voinov-1977,Cox-1986,Hocking-1992} %
\begin{align} \label{S7E4} %
  \dot s \ %
  = \ \frac{\ang^2}{\ln \frac 1\eps} (\ang-\gam),
\end{align}
which is valid with the assumption $n \in (1,3)$ in the limit $\eps \to
0$. Notice that the speed of the contact line vanishes if the outer contact
angle $\gam$ is the same as the microscopic contact angle $\ang$. The contact
line speed logarithmically converges to zero if $\eps \to 0$. Hence, in real
situations the change of the droplet profile $h$ and the motion of the contact
line $s(t)$ could take place in similar time scales. We note that the above
computation has some analogies with the computation in Subsection
\ref{MatchTypeb}.

\medskip

\paragraph{Complete wetting case.} Although most of the paper is concerned with
the case of a non--zero contact angle, for the convenience of the reader we also
give a short motivation of Tanner's Law which is relevant for the zero contact
angle case. We focus on the relevant case $n=2$, associated to the Navier slip
condition. We need to obtain a solution of \eqref{S5E9} with a suitable choice
of $\dot s$ and the boundary condition $H_y\left(0,\tau \right) =0$ such that
$h = \eps H$ satisfies the matching condition \eqref{S5E8}. We look for
solutions of \eqref{S5E9} in the form $H= H_{0}+H_1 + \ldots$ where
$|H_1| \ll H_{0}$ for $y \ll 1$.  The leading order contribution is determined
by $H_{0}H_{0,yyy} = \dot s$ (with $\dot s = \ddt s$) which yields
\begin{align}
H_0(y,\tau)  \  \approx \ \sqrt{\tfrac 83} (-\dot s)^{\frac 12}y^{\frac 32} \qquad\qquad %
  \text{as $y\to 0^+$}  \label{S6E3}
\end{align}
using a standard asymptotic argument.  The nonnegativity of the solution and
\eqref{S6E3} then show that $\dot s <0$, i.e. the support of the solutions can
only expand. In leading order \eqref{S5E9} then takes the form
$H_0^2 H_{0,yyy} + H_0 H_{1,yyy} + H_1 H_{0,yyy}$ = $0$ which leads to
\begin{align} \label{H1-nonhom} %
  \frac 83 y^3H_{1,yyy}-H_1 \ = \ \frac 83 (-\dot s) y^3.
\end{align}
The solution of \eqref{H1-nonhom} has the form $H_1 = H_{1,p} + H_{1,\hom }$
with the particular solution $H_{1,p} = \frac 8{45}(-\dot s) y^3$ and a solution
$H_{1,\hom}$ of the homogeneous problem. In particular, $H_{1,\hom}$ is a linear
combination of the functions $y^{\bet_i}$, $1 = 1,2,3$ with
$\bet_{1,2} = \frac{5}{4} \pm \frac 1{4}\sqrt{13}$ and $\bet_3 = \frac 12$. The
assumption $|H_1| \ll H_0$ implies $\bet_2 = \bet_3 = 0$. This yields for some
$K\in \R$ which might depend on $\dot s$ the asymptotic formula
\begin{align}
  H \ %
  \approx \ \sqrt{\frac{8}3} (-\dot s)^{\frac 12}y^{\frac 32}+ \frac 8{45} (-\dot s) y^3+Ky^{\bet_1}
  \qquad  %
  \text{for $(-\dot s)^{\frac 13} y \ll 1$.}
  \label{S6E4}
\end{align}
The asymptotic formula \eqref{S6E4} holds as long as $|H_1| \ll H_0$, i.e.
$y \ll (-\dot s)^{-\frac 13}$. Here, we make the self--consistent assumption
that $Ky^{\bet_1}$ yields a similar contribution than the other two terms for
these values of $y$. In terms of $\eta =(-\dot s)^{\frac 13}y$ this yields the
matching condition
\begin{align} \label{S6E6} %
  H \ \approx \ \sqrt{\tfrac 83} \, \eta^{\frac 32}+ \frac{8}{45}\eta^3+K(-\dot
  s)^{-\frac{\bet_1}3}\eta^{\bet_1} \qquad\qquad %
  \text{for $\eta \ll 1$.}
\end{align}
On the other hand, in terms of $\eta$, \eqref{S5E9} turns into %
\begin{align}
  1+ (H^2+H) H_{\eta \eta \eta } \ %
  = \ 0 \qquad\qquad %
  \text{for $\eta >0$},  \label{S6E5}
\end{align}%
Equation \eqref{S6E5} has been rigorously studied in
\cite{GiacomelliOtto-2002,GnannWisse-2022}. The behaviour of the solutions
of \eqref{S6E5} depends strongly on $K$: Except for a discrete set of values we
have either $H \approx y^2$ as $\eta \to \infty$ or $H =0$ for some positive
value of $\eta$. On the other hand, with the choice
$K=K_{0}\left(-\dot s \right)^{\frac{\bet_1}3}$ the solution satisfies
\begin{align}
  H \ %
  \approx  \ 3^{\frac 13}\eta (\ln \eta)^{\frac 13}\ \qquad\qquad %
  \text{as $\eta \to \infty$}.   \label{S6E7}
\end{align}
In terms of the initial variable $h$, cf. \eqref{inner-var-a}, equation \eqref{S6E7}
takes the form
\begin{align} \label{S7E1} %
  h \ %
  \approx \ 3^{\frac 13} (-\dot s)^{\frac 13}\xi \Big(\ln \frac {(-\dot
  s)^{\frac 13}\xi}{\eps} \Big)^{\frac 13} \ %
  \approx \ 3^{\frac 13} (-\dot s)^{\frac 13}\xi \big(\ln \frac 1\eps
  \big)^{\frac 13} \qquad %
  \text{for $\xi \sim 1$}
\end{align}
where we used the assumption $|\ln (-\dot s)^{\frac 13}\xi| \ll \ln \frac 1\eps$
for the last approximation which is consistent with \eqref{S7E2}.  Matching the
behaviour \eqref{S7E1} in the slip region with the asymptotics \eqref{S5E8} we
obtain
\begin{align}
  \dot s \ %
  = \ - \frac{\gam^3}{3 \ln \frac 1{\eps}} %
  \qquad\qquad \text{as $\eps \to 0$}.  \label{S7E2}
\end{align}
A similar argument can be worked out in the more general case $n < 3$ noting
that a logarithm appears on certain critical values of $n$.


\subsection{Approximating type (b) solutions \label{MatchTypeb}}

In this subsection we consider the limit $\eps \to 0$ for solutions of
\eqref{tfe-slip} in the case of type (b) solutions. We consider the situation
equation \eqref{tfe-slip-xi} of a single contact point where the solution has
support $(s(t), \infty)$. We assume that $n \in (0,3)$ is fixed.  At the contact
point we assume the microscopic contact angle
$\ang_\eps \approx \gam (-\ln \eps)^{\frac 13}$.  Furthermore, the solution of
\eqref{tfe-slip} approximately has the behaviour \eqref{ass-type-b} of a type
(b) solutions of \eqref{tfe-noslip} in the outer region for some $\OO(1)$
function $\gam$ in the outer region where $h^3\gg \eps^{3-n}h^{n}$
holds. 

\medskip

The transition between inner and outer regions occurs at the location where
$h^3 \approx \eps^{3-n}h^{n}$. Inserting the expansion \eqref{ass-type-b}, the
size $\Delta_\eps$ of the inner region is determined by
$\Delta_\eps^3 (- \ln \Delta_\eps) \ %
\approx \ \eps^{3-n}\Delta_\eps^n (-\ln \Delta_\eps)^{\frac{n}{3}}$,
% or equivalently $\Delta_\eps^{3-n}(- \ln \Delta_\eps )^{\frac{3-n}{3}}\approx \eps^{3-n}$,
i.e. $\Delta_\eps \ \approx \ \eps (-\ln \eps)^{-\frac 13}$ as $\eps \to 0$
using that $\ln \Delta_\eps \approx \ln \eps$. We introduce new
variables by
  \begin{align}     \label{inner-var-b} %
    \xi \ =: \ \eps (-\ln \eps)^{-\frac 13}y, \quad %
    t \ =: \ t_{0}+ \eps (-\ln \eps)^{-\frac 43}\tau,\quad %
    h \ =: \ \eps H,\  \ %
    % s(t) = \eps (-\ln \eps)^{-\frac 13}\sig(\tau),  %
  \end{align}
  where $t_0 > 0$ is a fixed time. With the assumption
  $\ang_\eps \approx \gam (-\ln \eps)^{\frac 13}$, $H$ solves
\begin{align}
  &H_{\tau }-(-\ln \eps)^{-1}%
  \dot sH_{y}+((H^3+H^{n})H_{yyy})_{y} \ = \ 0, \qquad %
  &&\text{for $y>0$},  \label{T4E2} \\
  &H \ = \ 0,\ \ H_{y} \ = \ \gam \ \ \ %
    &&\text{for $y=0$,}  \label{T4E3-2}
\end{align}
where as always $\dot s = \frac d{dt}s$. Since small changes in $t$ correspond
to huge changes in the new time variable $\tau$ we approximate 
\eqref{T4E2} by the steady state equation
\begin{align}
  &-(-\ln \eps)^{-1}\dot sH_{y}+((H^3+H^{n})H_{yyy})_{y} =0, \qquad %
    &&\text{for $y>0$.}  \label{T4E4}
\end{align}
Integrating \eqref{T4E4}  and dividing by $H$ we obtain%
\begin{align}
  -(-\ln \eps)^{-1}\dot s+(H^2+H^{n-1})H_{yyy} \ \lupref{no-massflux}= \ 0,%
  &&\text{for $y>0$}
  \label{T5E3}
\end{align}
To the leading order in $\eps \to 0$, $H$ hence satisfies
\begin{align}
  &(H^2+H^{n-1})H_{yyy} \ = \  0,\ \ &&\text{for $y > 0$},\ \\
  &H=0,\ \ H_{y} = \gam \ \ \ &&\text{for $y=0$.} \label{eq-H2} %
\end{align}
In view of \eqref{ass-type-b} for $\Delta_\eps \ll \xi \ll 1$, i.e.
$1 \ll y \ll \eps (-\ln \eps)^{-\frac 13}$, we obtain the matching condition
\begin{align} \label{T4E6} %
  H \ %
  \lupref{inner-var-b}= \ \frac h\eps \ %
  &\lupref{ass-type-b}\approx \ \frac{\gam \xi}\eps (- \ln \xi)^{\frac 13} \ %
    \lupref{inner-var-b}\approx \ \gam y \Big( 1 - \frac{\ln y}{(-\ln \eps)} \Big)^{\frac 13} \ %
    \approx \ \gam y - \frac{\gam y \ln y}{3(-\ln \eps)}, %
\end{align}%
where we used $-\ln \xi \approx -\ln (\eps y) \approx (-\ln \eps) - \ln y$.
Hence $H\approx \gam y$ for $1 \ll y \ll \eps (-\ln \eps)^{-\frac 13}$. Since
$H$ solves \eqref{T5E3}, this also implies $H\approx \gam y$ for $y \lesssim 1$.  To
get the next order correction we write
\begin{align}
  H  \ %
  =: \ \gam y+ \frac 1{(-\ln \eps)}U(y)  \qquad\qquad %
  \text{for $y > 0$.} \label{T5E2}
\end{align}%
Plugging \eqref{T5E2} into \eqref{T5E3} we obtain to the leading order %
\begin{align}
  U_{yyy} \ = \ \frac {\dot s}{\gam^2 y^2+ \gam^{n-1} y^{n-1}} \qquad\qquad % 
  \text{for $y > 0$}.
\end{align}%
By assumption, we have $U_{yy} \to 0$ for $y \to \infty$. Integrating we hence
get
\begin{align} \label{eq-Uyy} %
  U_{yy}(y) \ %
  \approx \ - \frac {\dot s}{\gam^2} Q_\gam(y), \qquad
  \text{where } Q_\gam(y) :=  \int_{y}^{\infty }\frac 1{z^2+ \gam^{n-3}z^{n-1}} \ dz \ %
\end{align}
noting that the integral is finite if $n < 3$. Furthermore, 
\begin{align} \label{lob-3} %
  \Big|Q_{\gam }(y) - \frac 1{y} \Big| \ %
  = \ \int_{y}^{\infty }\frac {\gam^{n-3}z^{n-1}}{(z^2+ \gam^{n-3}z^{n-1})z^2} \
  dz \ %
  % \approx \ \int_{y}^{\infty }\frac 1{z^2}dz= \frac 1{y} \quad %
  % \text{ as $y\to \infty$}.
  = \ o(\frac 1{y^2}) \qquad 
  \text{for  $y \gg 1$.}
\end{align}
Integrating \eqref{eq-Uyy} twice using $U(0) = U_y(0) = 0$ yields%
\begin{align} \label{lob-1} %
  U(y) \ %
  &= \ - \frac {\dot s}{\gam^2}\int_{0}^{y}\int_{0}^{z_1}Q_{\gam}(z_2) \ dz_2dz_1+ o(y \ln y) \\
  &= \ - \frac {\dot s}{\gam^2}y\ln y + o(y \ln y) %
  \qquad\qquad \text{for $y\gg 1$}.
\end{align}%
Matching \eqref{lob-1} with \eqref{T4E6} we obtain formula \eqref{dots-formula},
i.e.
\begin{align} \label{dots-formula2} %
  \dot s \ %
  = \   + \lim_{\eps \to 0} \frac {\ang_\eps^3}{3|\ln \eps|} \ %
  = \  + \frac \gam 3. \qquad %
\end{align}

\section{Conclusion} %

We analyze the evolution of thin liquid droplets in the lubrication
approximation with different slip conditions at the liquid--solid
interface. Motivated by the classical no-slip paradox which states that the
Navier–Stokes equations with a no-slip boundary condition require unphysical
infinite dissipation during droplet spreading, we focus on the limit of
vanishing slip.  Our analysis shows that there are three different physically
relevant limits where the solutions approximate qualitatively different
evolutions of the free boundary: We show that in the no--slip limit three
fundamentally different classes of limiting solutions are approached, each of
them corresponding to a different scaling of the microscopic contact angle as
the regularization parameter vanishes:
\begin{itemize}[leftmargin=5ex]
\item[(a)] The first solution class are regular at the contact line and have
  fixed support. This reflects the well known fact that the contact line remains
  fixed for sufficiently regular no-slip solutions.
\item[(b)] For the second solution class, the contact line recedes on a time
  scale of order one while the solution's profile exhibits a characteristic
  logarithmic correction near the moving contact point. Intriguingly, for these
  solutions the dissipation near the contact point diverges, but this divergence
  is balanced by the energetic cost associated with the shrinking solid–liquid
  interface.
\item[(c)] A third class of solutions appears when the contact angle diverges
  even faster in which case the contact line recedes on a much smaller time
  scale, and the mass removed from the support accumulates at the boundary.
\end{itemize}
These findings suggest that the thin-film equation with no slip supports a rich
family of physically admissible solutions, provided one interprets the no--slip
thin film equation as the asymptotic limit of models which regularized slip
conditions. Even though the large apparent contact angles in solutions of type
(b) and (c) seem incompatible with the lubrication approximation, a refined
analysis shows that the underlying physical variables remain consistent with the
assumptions for the lubrication approximation.


\bibliographystyle{plain} %
\bibliography{tfe}

\end{document}
